\documentclass[11pt]{article}
\usepackage[margin=1in]{geometry}
\usepackage{amsmath,amssymb,amsthm,mathtools,array,booktabs,microtype}
\usepackage[colorlinks=true,linkcolor=blue,citecolor=blue,urlcolor=blue]{hyperref}
\hypersetup{pdftitle={Regularized Scaling for Lossy Generalized Flow},
pdfsubject={Clarified proof draft with expanded derivations},pdfauthor={Research manuscript}}
\newtheorem{theorem}{Theorem}[section]
\newtheorem{lemma}[theorem]{Lemma}
\newtheorem{proposition}[theorem]{Proposition}
\newtheorem{corollary}[theorem]{Corollary}
\theoremstyle{definition}\newtheorem{definition}[theorem]{Definition}
\theoremstyle{remark}\newtheorem{remark}[theorem]{Remark}
\newcommand{\R}{\mathbb R}
\newcommand{\one}{\mathbf1}
\newcommand{\diag}{\operatorname{diag}}
\newcommand{\norm}[1]{\left\lVert#1\right\rVert}
\newcommand{\poly}{\operatorname{poly}}
\DeclareMathOperator{\OPT}{OPT}
\setlength{\emergencystretch}{2em}
\setlength{\parskip}{2pt}
\title{Regularized Scaling for Lossy Generalized Flow\\
\large A clarified unconditional proof draft}
\author{Research manuscript}
\date{22 October 2026}
\begin{document}
\maketitle
\begin{abstract}
The proposed algorithm optimizes a dual vector and a positive vertex
scaling together. Its potential is designed to produce two certificates:
an ordinary grounded graph approximates the scaled generalized incidence
norm, and the primal imbalance admits a small weighted correction.
The correction gives an exactly feasible nonnegative flow in a padded
linear program. Three scalar identities bound its primal--dual gap.

The implementation follows the same structure. One local derivative
bound controls both curvature and the error of cached complete rows.
Ordinary accumulated-motion detectors identify which rows to refresh,
including rows internal to a scaling cut. Potential decrease pays for
these refreshes and for occasional anchor resets. The main text follows
this proof chain; the appendices give range estimates, the original-flow
reduction, exact maximum-delivery repair, and numerical details.
\end{abstract}
\noindent\textbf{Proof status.} This is a proposed proof of the stated
unconditional bound, with the arithmetic model and output guarantees
specified below. The new arguments remain subject to independent
mathematical verification.
\setcounter{tocdepth}{1}
{\small\tableofcontents}
\newpage

\section{The claim and the proof in one page}

For an arc $e=(u,v)$, sending $f_e$ units consumes $f_e$ at $u$ and
supplies $\gamma_e f_e$ at $v$. Let
$B_e=-e_u^{\mathsf T}+\gamma_e e_v^{\mathsf T}$, so
$B^{\mathsf T}f$ is net inflow. The source supplies arbitrary flow;
maximum delivery means maximizing net inflow into the sink.

Let $M=m_{\rm original}+n_{\rm original}+2$, let $0<\delta<1$, and
let $W\ge2$ bound original costs, capacities, demands, gains, and the
inverses of nonzero original data. Put
\[
 H=100+\lceil\log_2(MW/\delta)\rceil.
\]
After the linear-size reduction, $m,n$ denote the reduced numbers of
rows and coordinates; both are $O(M)$. No lower bound is assumed for
quantities formed by cancellation, such as a reduced demand or a gradient.

\begin{theorem}[Proposed unconditional bound]\label{thm:main}
For a feasible generalized minimum-cost flow instance, the algorithm
returns $0\le\bar f\le u$ with
\[
 \norm{B^{\mathsf T}\bar f-d}_\infty\le\delta,
 \qquad c^{\mathsf T}\bar f\le\OPT+\delta.
\]
For lossy maximum delivery, $0<\gamma_e\le1$, it returns a capacitated
flow with exact internal conservation and delivery at least
$\OPT-\delta$. The arithmetic work is
$M^{1+o(1)}\poly(H)$, with high probability.
\end{theorem}

The output distinctions are essential:
\begin{center}
\begin{tabular}{p{.29\linewidth}p{.61\linewidth}}
\toprule
Problem & Feasibility guarantee\\\midrule
Padded linear program & Exact equality and nonnegativity.\\
Original minimum-cost flow & Exact capacities; additive demand error at most $\delta$.\\
Lossy maximum delivery & Exact capacities and internal conservation, after a separate repair.\\
\bottomrule
\end{tabular}
\end{center}
The optimization uses $\poly(H)$ working bits. Exact residual corrections
and final flow sweeps use original-coefficient arithmetic and shared
expressions. The operation bound does not bound the expanded encoding
length of arbitrary rational inputs or products of their gains.

One phase starts with a feasible flow $f^0$ and a feasible dual vector
$a$. It reduces their certified gap by a constant factor. The logical
order of the proof is:
\begin{enumerate}
\item Identify the correction that would produce a better feasible flow.
\item Construct a potential whose stationary points certify that correction
and bound the recovered gap.
\item Prove that complete rows can be kept current using ordinary graph
motion detectors.
\item Charge the cut steps, refreshes, and anchor resets to potential decrease.
\end{enumerate}
The reverse norm comparison is required only at the end of a phase.
It is not an invariant assumed during descent. The scaling itself can
change on a large cut; the amortized recourse bound concerns its cached
row representation.

\section{What one phase must produce}\label{sec:phase-target}

\subsection{The grounded linear program and phase data}

The reduced primal--dual pair is
\begin{equation}\label{eq:LP}
 \min\{c^{\mathsf T}f:A^{\mathsf T}f=d,\ f\ge0\},
 \qquad \max\{d^{\mathsf T}y:Ay\le c\}.
\end{equation}
Each row has at most two nonzeros, with opposite signs if there are two.
Every coordinate also has the signed unit unary rows $+e_i^{\mathsf T}$
and $-e_i^{\mathsf T}$, with positive costs. These ground the dual
coordinates and provide an elementary exact residual correction.
Appendix~\ref{app:reduction} gives the reduction and explains what is
lost when artificial variables are removed.

Fix a nonnegative feasible baseline $f^0$ and a strictly feasible phase
anchor $a_{\rm start}$. Let
\begin{equation}\label{eq:phase-data}
 F=c^{\mathsf T}f^0,\quad
 \Delta_{\rm start}=F-d^{\mathsf T}a_{\rm start}>0,\quad
 \lambda=K/\Delta_{\rm start},\quad w=\lambda f^0.
\end{equation}
Then $A^{\mathsf T}w=\lambda d$ and
$w^{\mathsf T}(c-Aa_{\rm start})=K$. Use an internal stopping tolerance
at most $\delta/100$; the remaining accuracy margin pays for rounding,
padding, clipping, and final repair.

The phase has a movable displacement anchor $a$, initially
$a_{\rm start}$, and two optimization fields
\begin{equation}\label{eq:coordinates}
 s_i=\log(q_i/a_0),\qquad x_i=(y_i-a_i)/q_i,
 \qquad y=a+\diag(q)x,
 \qquad \sigma=c-Ay>0.
\end{equation}
This centered logarithmic field starts at zero when $q=a_0\one$.
It has exactly the same derivatives and cut increments as $\log q$.
For a two-ended row, its owner $u_e$ is its positive endpoint:
$A_e=a_e e_{u_e}^{\mathsf T}-b_e e_{v_e}^{\mathsf T}$ with
$a_e,b_e>0$. Unary rows retain their actual sign.

Throughout, the Huber function and its derivative are
\[
 \phi(t)=\begin{cases}t^2/2,&|t|\le1,\\ |t|-1/2,&|t|>1,\end{cases}
 \qquad\theta(t)=\max\{-1,\min\{1,t\}\}.
\]

\begin{center}
\begin{tabular}{p{.27\linewidth}p{.64\linewidth}}
\toprule
Fixed within a phase & $f^0,F,\lambda,w,K$ and the barrier references
$\widehat S=c-Aa_{\rm start}$.\\
Updated by cut steps & $s,x$, and therefore $q,y,\sigma$ and all row weights.\\
Changed only by a reset & The displacement anchor $a$; its slack is
$S^a=c-Aa$ and its mass is $K_a=w^{\mathsf T}S^a$.\\
\bottomrule
\end{tabular}
\end{center}
In particular $K_a$ need not equal $K$. Keeping these quantities distinct
prevents a hidden change of the barrier when the anchor is reset.

\subsection{The recovery identity, before the potential}

Suppose a later computation supplies positive row values $\ell_e$.
Set $L=\ell+w$. We want a correction $v$ such that
\begin{equation}\label{eq:target-correction}
 A^{\mathsf T}(\ell-L\odot v)=\lambda d,
 \qquad \norm v_\infty\le\varepsilon<1.
\end{equation}
The flow $(\ell-L\odot v)/\lambda$ has exact conservation, but can
have negative entries. The baseline fixes that difficulty.

\begin{proposition}[Nonnegative recovery]\label{prop:early-recovery}
If \eqref{eq:target-correction} holds, then
\begin{equation}\label{eq:recovery-flow-formula}
 \boxed{f=\frac{(1-v)\odot\ell+(\varepsilon-v)\odot w}
                         {\lambda(1+\varepsilon)}}
\end{equation}
is nonnegative and satisfies $A^{\mathsf T}f=d$. For any feasible dual
$y$, let $T=\ell^{\mathsf T}\sigma$ and
$W_\sigma=w^{\mathsf T}\sigma$. Then
\begin{equation}\label{eq:recovery-gap-formula}
 c^{\mathsf T}f-d^{\mathsf T}y
 \le [T+2\varepsilon W_\sigma]/\lambda.
\end{equation}
\end{proposition}
\begin{proof}
Both coefficients in the numerator are nonnegative. That numerator is
$\ell-L\odot v+\varepsilon w$, whose image under $A^{\mathsf T}$ is
$(1+\varepsilon)\lambda d$. Finally the gap equals
$\sigma^{\mathsf T}f$. Its numerator is at most
$(1+\varepsilon)T+2\varepsilon W_\sigma$, proving the bound.
\end{proof}

Thus the potential has a concrete job: certify a small correction and
make $T$ small compared with $K$. It does not need to make $q$ change
in only a few coordinates.

% Replacement scalar-potential and range sections.  The preceding setup
% supplies phase data, sigma,x,s,phi,theta, and signed physical rows A.
\section{The potential, defined by one scalar equation}
\label{sec:potential}\label{sec:envelope-ranges}

The barrier reference and the displacement anchor have different jobs.
At phase start, fix
\[
 \widehat S_e=c_e-A_ea_{\rm start}>0.
\]
An anchor reset changes $a$ and hence the current anchor slacks
$S_e^a=c_e-A_ea$; it does not change $\widehat S_e$, $w$, or $\lambda$.
This distinction will matter in the phase-progress calculation.

Fix $B=64$, $\eta=1/8$, and a barrier exponent $0<\alpha_{\rm b}\le1$.
For $z>0$, define
\[
 \beta_e(z)=-B\log\frac z{S_{\max}}
       +\frac B{\alpha_{\rm b}}
              \left(\frac{\widehat S_e}{z}\right)^{\alpha_{\rm b}}.
\]
Here $S_{\max}$ is an upper bound on every feasible dual slack.
The logarithm keeps a slack positive.  The power term will turn a
potential bound into a bound on complementarity.

For the current $(y,q)$, put
\[
 r_e^0=q_{u_e}^{-\kappa}A_eq,\qquad
 t_e^0=\frac{q_{u_e}^{-\kappa}}{b_0}A_e(y-a),\qquad
 b_e=|r_e^0|+\eta|t_e^0|.
\]
The superscript $0$ means that these quantities do not contain the
unknown row weight.  For a trial $\ell>0$, let $z_e(\ell)>0$ be the
unique solution of
\begin{equation}\label{eq:barrier-inverse}
 \ell=-\beta'_e(z_e(\ell))
 =\frac B{z_e(\ell)}
       \left[1+\left(\frac{\widehat S_e}{z_e(\ell)}\right)
                         ^{\alpha_{\rm b}}\right].
\end{equation}
This inverse is well defined because $-\beta'_e$ decreases continuously
from $+\infty$ to zero.

The row computation is the following single scalar root:
\begin{equation}\label{eq:envelope-root-exact}
 \sigma_e=z_e(\ell_e)
   +r_e^0\theta((\ell_e+w_e)r_e^0)
   +\eta t_e^0\theta((\ell_e+w_e)t_e^0).
\end{equation}
After finding it, set
\[
 L_e=\ell_e+w_e,\quad z_e=z_e(\ell_e),\quad
 C=\diag(L_eq_{u_e}^{-\kappa})A\diag(q),\quad
 r=C\one,\quad t=Cx/b_0.
\]
Thus $r_e=L_er_e^0$ and $t_e=L_et_e^0$.  Define the row potential by
\begin{equation}\label{eq:row-root-value}
 \boxed{\quad
 \Psi_e=\beta_e(z_e)-\ell_e(\sigma_e-z_e)
                         +\phi(r_e)+\eta\phi(t_e).\quad}
\end{equation}
If the root does not exist, set $\Psi_e=+\infty$.

The full potential is
\begin{equation}\label{eq:potential}
 \Phi_a(y,q)=\lambda(F-d^\top y)+\sum_e\Psi_e+R(q)+E(x).
\end{equation}
Here
\[
 R(q)=\sum_i\omega_i\left[\frac{a_0}{q_i}-\log\frac{q_i}{a_0}\right],
 \qquad E(x)=\sum_i\frac{\omega_i}{p}(x_i/b_0)^p,
\]
and $\omega_i=64\deg_A(i)+1$.
The integer $p\ge2$ is even.  The inverse term in $R$ prevents a
scaling coordinate from approaching zero within a bounded sublevel;
the paired unary rows will bound it from above.  Neither conclusion
comes from the negative logarithm alone.

\begin{lemma}[The root and the exact finite domain]
\label{lem:envelope-root}\label{lem:envelope-domain}
The root \eqref{eq:envelope-root-exact} exists uniquely precisely when
\[
 \boxed{\sigma_e>|r_e^0|+\eta|t_e^0|.}
\]
On this domain the row potential is continuously differentiable and
its derivatives can be evaluated holding $\ell_e$ fixed.  With
$c_\eta=(1+\eta)/2$, it satisfies
\begin{align}
 \beta_e(\sigma_e-b_e)+w_eb_e-c_\eta
 &\le\Psi_e\le\beta_e(\sigma_e-b_e)+w_eb_e,
 \label{eq:scalar-sandwich}\\
 \Psi_e&\ge\beta_e(\sigma_e),
 \label{eq:barrier-lower}\\
 \sigma_e-b_e&\le z_e\le\sigma_e,
 \label{eq:root-slack-range}\\
 \Psi_e&\ge\beta_e(z_e)-c_\eta.
 \label{eq:root-value-lower}
\end{align}
\end{lemma}

\begin{proof}
Suppress the row subscript.  Write $z=z(\ell)$ and consider
\[
 J(\ell)=\beta(z)+\ell(z-\sigma)
       +\phi((\ell+w)r^0)+\eta\phi((\ell+w)t^0).
\]
The identity $\beta'(z)=-\ell$ cancels the derivatives of $z$, giving
\[
 J'(\ell)=z-\sigma+r^0\theta((\ell+w)r^0)
                         +\eta t^0\theta((\ell+w)t^0).
\]
For $v=(\widehat S/z)^{\alpha_{\rm b}}$,
\[
 \frac{\beta''(z)}{\ell^2}
   =\frac{1+(1+\alpha_{\rm b})v}{B(1+v)^2}
   \le\frac{1+\alpha_{\rm b}}B.
\]
Hence $z'(\ell)=-1/\beta''(z)\le-B/[(1+\alpha_{\rm b})\ell^2]$.
A Huber term has nonzero second derivative only when the absolute
value of its argument is below one.  Since $\ell+w\ge\ell$, the
two positive second derivatives sum to at most $(1+\eta)/\ell^2$.
Therefore, away from the Huber breakpoints,
\[
 J''(\ell)\le-
 \frac{B/(1+\alpha_{\rm b})-1-\eta}{\ell^2}<0.
\]
The first derivative is continuous and strictly decreasing.  Its limits
at zero and infinity are $+\infty$ and $-\sigma+b$, respectively.
It consequently has exactly one zero when $\sigma>b$ and none otherwise.
The zero is the unique maximum of $J$, proving that the root-defined
value is equivalently $\Psi=\sup_{\ell>0}J(\ell)$.

For $s>0$, the function $\beta(z(\ell))+\ell(z(\ell)-s)$ has
derivative $z(\ell)-s$ and maximum $\beta(s)$.  Together with
$|u|-1/2\le\phi(u)\le|u|$, this proves the sandwich.
If $\sigma\le b$, the lower Huber bound instead gives
\[
 J(\ell)\ge\beta(z(\ell))+\ell z(\ell)
                    +\ell(b-\sigma)+wb-c_\eta\longrightarrow+\infty.
\]
Indeed $z(\ell)\to0$ and $\beta(z(\ell))\to+\infty$.
This also proves the exact finite-domain assertion from the variational
description.  Choosing the trial value $\ell=-\beta'(\sigma)$ and
deleting the nonnegative Huber terms proves \eqref{eq:barrier-lower}.

Every product $r^0\theta((\ell+w)r^0)$ is between zero and $|r^0|$;
the same holds for $t^0$.  The root equation therefore gives
\eqref{eq:root-slack-range}.  At the root, put
$F_r=r\theta(r)+\eta t\theta(t)$.  Then
\[
 \Psi=\beta(z)-\frac\ell L F_r+\phi(r)+\eta\phi(t)
       \ge\beta(z)-c_\eta,
\]
because $0<\ell/L\le1$ and $\phi(u)-u\theta(u)\ge-1/2$.
Finally, the strict local derivative bound keeps the unique maximizer
in a compact interval under a small parameter perturbation.  Its
continuous dependence and the elementary envelope theorem give the
stated differentiability rule.
\end{proof}

\begin{remark}[Relation to the conjugate formula]
If $\beta_e^*$ denotes the convex conjugate, then
$-\beta_e^*(-\ell)=\beta_e(z_e(\ell))+\ell z_e(\ell)$.
Thus the variational formula just proved is exactly the conjugate
envelope.  The algorithm only needs the scalar equation
\eqref{eq:envelope-root-exact}.
\end{remark}



\section{The derivatives and the certificates they are designed to give}
\label{sec:gradients}

Three gradients occur because a scaling update with $x$ fixed also
changes the physical dual vector $y$. Define
\[
 g=\diag(q)A^{\mathsf T}(\ell-w),\quad
 e_i=\frac{\omega_i}{b_0}(x_i/b_0)^{p-1},\quad
 v_i=\omega_i[1+a_0/q_i+(x_i/b_0)^p],
\]
\[
 h_i=v_i+\kappa\sum_{e:u_e=i}
              [r_e\theta(r_e)+\eta t_e\theta(t_e)].
\]
All terms in $h$ are nonnegative, and $h_i\ge\omega_i$.
The derivatives are
\begin{align}\label{eq:gradients}
 g_q&=C^{\mathsf T}\theta(r)-h,
 &&\text{logarithmic scaling derivative at fixed physical }y,\\
 g_x&=g+\frac\eta{b_0}C^{\mathsf T}\theta(t)+e,
 &&\text{derivative with respect to }x\text{ at fixed }s,\nonumber\\
 G&=g_q+\diag(x)g_x,
 &&\text{derivative with respect to }s\text{ at fixed }x.\nonumber
\end{align}
The algorithm optimizes $G$ and $g_x$. The scaling certificate uses
$g_q$. Section~\ref{sec:recovery-endpoint} proves the norm estimate
needed to pass from the former pair to the latter.

\subsection{Deriving the formulas}

The root construction permits envelope differentiation: hold $\ell_e$
fixed when taking a first derivative. For an $x$-variation,
$dy=D_qdx$, and
\[
 d\sigma=-AD_qdx,\qquad
 L_e\,dt_e^0=C_e dx/b_0,\qquad d r_e^0=0.
\]
The linear objective contributes $-A^{\mathsf T}w$, the slack term
contributes $A^{\mathsf T}\ell$, the displacement Huber term contributes
$(\eta/b_0)C^{\mathsf T}\theta(t)$, and $E$ contributes $e$.
This proves the formula for $g_x$.

At fixed physical $y$, $dx=-x\odot ds$. The row variations are
\[
 L_e,dr_e^0=C_e ds-\kappa r_e ds_{u_e},\qquad
 L_e,dt_e^0=-\kappa t_e ds_{u_e}.
\]
The vertex regularizers contribute
$-\omega_i(1+a_0/q_i)-\omega_i(x_i/b_0)^p$.
Combining these terms gives $g_q$. Finally, at fixed $x$,
$dy=D_q\diag(x)ds$, so the chain rule gives
$G=g_q+\diag(x)g_x$. The apparent scaling derivative of $E$ cancels
in $G$, because $E(x)$ is independent of $s$ in the actual coordinates.

\subsection{The only range facts needed in the main argument}

The unary rows and the potential give explicit ranges; their proofs are
in Appendix~\ref{sec:paired-unary-ranges}. On a phase sublevel
$\Phi\le P$, write
\[
 D=P+\Omega\log(q_{\max}/a_0),\qquad
 \Omega=\sum_i\omega_i,\qquad q_{\max}=C_0^{1/(1-\kappa)},
\]
where $C_0\ge\max|c_e|$ also bounds the paired unary costs.
Then
\begin{equation}\label{eq:range-summary}
 \lambda(F-d^{\mathsf T}y)+\sum_e\Psi_e+
 \sum_i\omega_i a_0/q_i+E(x)\le D,
 \quad q_i\ge\omega_i a_0/D,
 \quad |x_i|\le b_0(pD/\omega_i)^{1/p}.
\end{equation}
Moreover the finite-domain margin is positive, all scalar row weights
have logarithmic range $\poly(H)$, and
\begin{equation}\label{eq:Z-summary}
 Z:=\sum_e\ell_ez_e\le Bm+\alpha_{\rm b}D+
                         \alpha_{\rm b}m(1+\eta)/2.
\end{equation}
These facts justify the local derivatives before stationarity is known.
An additional endpoint argument, using the norm certificate, proves
$q_i>a_0$ and hence $\sum_i\omega_i(1+a_0/q_i)\le2\Omega$.
It is used only after the endpoint test; it is not assumed during descent.

% Replacement endpoint/recovery section. The preceding text defines the LP,
% potential, r=C1, t=Cx/b_0, and all three gradients in eq:gradients.
% Requires amsmath, amssymb, amsthm. Citation key: sherman.

\section{What stationarity gives, and how to recover the flow}
\label{sec:recovery-endpoint}

The endpoint argument has a short logical chain. The two coordinate
gradients certify a good ordinary graph norm. They also certify that the
primal imbalance has a small correction in the required weighted norm.
A general minimum-norm solver constructs that correction. Finally, three
scalar identities bound the primal--dual gap of the recovered flow.
No dynamic invariant is used in these static arguments.

\subsection{The scaling certificate in one absorption argument}

For a two-ended row $C_e=Q_e e_u^{\mathsf T}-R_e e_v^{\mathsf T}$,
$Q_e,R_e\ge0$, use the ordinary row norm
\[
 \rho_e(z)=R_e|z_u-z_v|+|Q_e-R_e||z_u|.
\]
For $\|b\|_{N,*}=\sup_{z\ne0}|b^{\mathsf T}z|/N(z)$ we use the usual dual norm. A unary row contributes its absolute coefficient times $|z_u|$; its
physical sign is retained in $C$. Write
\[
 \rho=\sum_e\rho_e,\qquad N(z)=\rho(z)+h^{\mathsf T}|z|.
\]
Thus $N$ is an ordinary graph norm with positive ground edges. Always,
\begin{equation}\label{eq:recovery-proxy-basic}
 \|Cz\|_1\le\rho(z)\le N(z),\qquad N(|z|)\le N(z).
\end{equation}
The reverse comparison will follow from stationarity; it is not assumed.

\begin{lemma}[Scaling and displacement certificates]
\label{lem:recovery-absorption}
Suppose $\|x\|_\infty\le X$, $0<\kappa,\eta\le1$, and put
\[
 L_x=\frac{2(X+b_0/\eta)}\kappa,\qquad \Lambda=8/\kappa.
\]
Then $N(x\odot z)\le L_xN(z)$ for every $z$. If
\[
 \|G\|_{N,*}\le\tau_s,\qquad \|g_x\|_{N,*}\le\tau_x,
 \qquad \delta_s:=\tau_s+L_x\tau_x\le\kappa/32,
\]
then
\begin{equation}\label{eq:recovery-norm-certificate}
 \|g_q\|_{N,*}\le\delta_s,\qquad
 \|Cz\|_1\le N(z)\le\Lambda\|Cz\|_1,\qquad
 h^{\mathsf T}|z|\le2\|Cz\|_1.
\end{equation}
\end{lemma}

\begin{proof}
First, multiplication by $x$ has a controlled cost. On a two-ended row,
\[
 \rho_e(xz)\le X\rho_e(z)+R_e|x_u-x_v||z_u|,
 \qquad R_e|x_u-x_v|\le b_0|t_e|+X|r_e|.
\]
For the Huber derivative,
\begin{equation}\label{eq:recovery-huber-elementary}
 |a|\le a\theta(a)+\tfrac14.
\end{equation}
Consequently, the sum of the extra coefficients at owner $i$ is bounded
by the energy terms $\kappa\sum_{e:u_e=i}(r_e\theta(r_e)+
\eta t_e\theta(t_e))$ in $h_i$, plus
$(b_0+X)\deg_A(i)/4$. The latter is paid by
$\omega_i\ge64\deg_A(i)$. This gives
\[
 \rho(xz)\le X\rho(z)+
 \max\{X/\kappa,b_0/(\kappa\eta),(b_0+X)/256\}\,h^{\mathsf T}|z|.
\]
Adding $h^{\mathsf T}|xz|\le Xh^{\mathsf T}|z|$ proves the claimed
$L_x$ bound. Unary rows need only the direct $X$ bound.

The gradient identity $g_q=G-\operatorname{diag}(x)g_x$ now gives
$\|g_q\|_{N,*}\le\delta_s$. The rest is the following absorption.
For $D_z=\|Cz\|_1$ and $H_z=h^{\mathsf T}|z|$,
\begin{align*}
 H_z&\le D_z+\delta_sN(z),
       &&\text{test }g_q=C^{\mathsf T}\theta(r)-h\text{ against }|z|,\\
 \rho(z)&\le D_z+(2/\kappa)H_z,
       &&\text{use }\rho_e(z)\le|C_ez|+2|r_e||z_u|,\\
 N(z)&\le D_z+(1+2/\kappa)H_z,\\
 [1-(1+2/\kappa)\delta_s]N(z)&\le(2+2/\kappa)D_z,\\
 N(z)&\le(8/\kappa)D_z,\qquad H_z\le2D_z.
\end{align*}
In the first line, $\|C|z|\|_1\le\|Cz\|_1$ for opposite-sign
two-ended rows, and $N(|z|)\le N(z)$. In the second line,
\eqref{eq:recovery-huber-elementary} gives
$\sum_{e:u_e=i}|r_e|\le h_i/\kappa$. The final line uses
$\delta_s\le\kappa/32$.
\end{proof}

This proves the endpoint certificate directly from the actual coordinate
gradients $G,g_x$. There is no need to assume that the fixed-$y$ gradient
$g_q$ was separately optimized, or that the reverse norm comparison held
earlier in the phase.

\subsection{The correct norm for primal recovery}

Define
\[
 B_0=\operatorname{diag}(L)A\operatorname{diag}(q),\qquad
 0<q_{\kappa,-}\le q_{u_e}^{\kappa}\le q_{\kappa,+}\quad\text{for every row }e.
\]
We need $B_0^{\mathsf T}v=g$ with small $\|v\|_\infty$.
This is a weighted correction for $C$, because
$B_0=\operatorname{diag}(q_{u_e}^{\kappa})C$.
Write $\operatorname{opt}_{B_0}(b)=
\min\{\|v\|_\infty:B_0^{\mathsf T}v=b\}$.

\begin{lemma}[A weighted correction certificate]
\label{lem:recovery-radius}
A vector $b$ has a correction of radius $\varepsilon$ if and only if
\begin{equation}\label{eq:recovery-zonotope}
 |b^{\mathsf T}z|\le\varepsilon\|B_0z\|_1\quad\text{for every }z.
\end{equation}
Under Lemma~\ref{lem:recovery-absorption},
\begin{equation}\label{eq:recovery-radius}
 \operatorname{opt}_{B_0}(g)
 \le\frac{\Lambda\tau_x+(\eta+2)/b_0}{q_{\kappa,-}}.
\end{equation}
\end{lemma}

\begin{proof}
The compact convex set $B_0^{\mathsf T}[-\varepsilon,\varepsilon]^m$
has support function $\varepsilon\|B_0z\|_1$. Separation proves
\eqref{eq:recovery-zonotope}, including sufficiency.
For the radius bound, use the displacement-gradient identity in
\eqref{eq:gradients} and
$|e_i|\le h_i/b_0$ to get
\[
 |g^{\mathsf T}z|
 \le\tau_xN(z)+(\eta/b_0)\|Cz\|_1+b_0^{-1}h^{\mathsf T}|z|
 \le[\Lambda\tau_x+(\eta+2)/b_0]\|Cz\|_1.
\]
Now $\|B_0z\|_1\ge q_{\kappa,-}\|Cz\|_1$, proving
\eqref{eq:recovery-radius}.
\end{proof}

The correction equation $B_0^{\mathsf T}v=g$ is equivalent to
$A^{\mathsf T}(\ell-L\odot v)=\lambda d$ after cancelling $D_q$.
Proposition~\ref{prop:early-recovery} therefore gives the nonnegative
flow and its gap bound, with no further feasibility assumption.

In particular, a bound on an unweighted solution of $C^{\mathsf T}z=g$
cannot replace \eqref{eq:recovery-radius}. The factor
$q_{\kappa,-}^{-1}$ records precisely the additional range requirement.

\subsection{Constructing a correction of a known radius}

The radius certificate gives more information than the recovery algorithm
needs: an explicit positive number
\[
 r_*:=\frac{\Lambda\tau_x+(\eta+2)/b_0}{q_{\kappa,-}}
 \quad\text{such that}\quad \operatorname{opt}_{B_0}(g)\le r_*.
\]
We construct a correction of norm at most $3r_*/2$. There is no search for
the possibly much smaller optimum and no precision requirement relative
to that optimum. This is useful when cancellation makes the true
imbalance arbitrarily small, or exactly zero.

The imported object is an ordinary \emph{congestion approximator}:
\begin{equation}\label{eq:recovery-preconditioner-import}
 \|Rb\|_\infty\le\operatorname{opt}_{D_N}(b)
                 \le\chi\|Rb\|_\infty
 \qquad\text{for every nonground demand }b.
\end{equation}
Here $N(z)=\|D_Nz\|_1$. Sherman's Theorem~6.1~\cite{sherman}, with
the normalized cut-congestion construction in his Section~6.1, builds
such a linear map with $\chi=n^{o(1)}$ in $M^{1+o(1)}$ work and applies
$R,R^{\mathsf T}$ in $n^{1+o(1)}$ work. Ground demands are included by
augmenting $b$ with $-\sum_i b_i$. Each normalized cut row has the form
$\mathbf1_S^{\mathsf T}b/\operatorname{cap}(\partial S)$, which proves
the first inequality. The second is the congestion approximation
property. The import explicitly includes both inequalities; it is not
merely a condition bound for a potentially noninjective map. In
particular it makes $R$ injective on the demand space. This is the
left-cancellable preconditioner convention stated in Section~5 of the
same paper.

\begin{proposition}[Known-radius correction]
\label{thm:recovery-constructive}
Set
\[
 \kappa_0=\chi\Lambda q_{\kappa,+}/q_{\kappa,-},\qquad
 M_0=\max_i\frac{\sum_eL_e|A_{ei}|}{L_{\rm unary(i)}}
       \le mA_{\max}L_{\max}/L_{\min},\qquad A_{\max}=\max_{e,i}|A_{ei}|.
\]
Given
$\operatorname{opt}_{B_0}(g)\le r_*$, one can construct an exact
correction of norm at most $3r_*/2$ using one ordinary preconditioner
build and
\begin{equation}\label{eq:recovery-static-iterations}
 O\bigl(\kappa_0^2\log(2k)[1+\log M_0]\bigr)
\end{equation}
products by $B_0,B_0^{\mathsf T},R,R^{\mathsf T}$ and simple vector
operations. Here $k$ is the number of rows of $R$.
Every product costs $M^{1+o(1)}$ work. Thus the static solve costs
$M^{1+o(1)}\operatorname{poly}(H)$ under the previously proved range
bounds. Its smallest stage radius is at least $r_*/(24M_0)$.
\end{proposition}

\begin{proof}
First dualize the endpoint norm comparisons and use
\eqref{eq:recovery-preconditioner-import}. For every demand $b$,
\begin{equation}\label{eq:recovery-preconditioned-bounds}
 \|Rb\|_\infty\le q_{\kappa,+}\operatorname{opt}_{B_0}(b),
 \qquad
 \operatorname{opt}_{B_0}(b)\le
      \frac{\chi\Lambda}{q_{\kappa,-}}\|Rb\|_\infty.
\end{equation}
We describe a constant-factor residual reduction, then repeat it.

Suppose a residual $h$ satisfies $\operatorname{opt}_{B_0}(h)\le\rho$.
For $\|u\|_\infty\le\rho$, define $2k$ signed losses
\[
 \mathcal L_{a,\pm}(u)=
 \frac{\pm[R(B_0^{\mathsf T}u-h)]_a}{q_{\kappa,+}\rho},
 \qquad a=1,\ldots,k.
\]
Both terms in the numerator have norm at most $q_{\kappa,+}\rho$,
by \eqref{eq:recovery-preconditioned-bounds}; therefore all losses lie
in $[-2,2]$. There is a feasible witness $u^*$ in this cube with every
loss zero. For any probability vector $p_t$ on the signed coordinates,
we can consequently choose a best response $u_t$ with
\[
 p_t^{\mathsf T}\mathcal L(u_t)\le0.
\]
This best response is explicit. If $p_t^+,p_t^-$ are its two halves,
\[
 u_t=-\rho\,\operatorname{sign}
       \bigl(B_0R^{\mathsf T}(p_t^+-p_t^-)\bigr).
\]
The convention at zero can be arbitrary.

Here is the required multiplicative-weights estimate, including its
constants. Start all $2k$ weights at one, normalize them to $p_t$, and
multiply weight $a$ by
$\exp(\vartheta\mathcal L_a(u_t))$, where
$\vartheta=1/(64\kappa_0)$. Since $|\vartheta\mathcal L_a|\le1$,
$e^z\le1+z+z^2$ gives an upper bound
\[
 \log Z_{T+1}\le\log(2k)
       +\vartheta\sum_{t=1}^T p_t^{\mathsf T}\mathcal L(u_t)
       +4\vartheta^2T
 \le\log(2k)+4\vartheta^2T
\]
for the total weight. The weight of an individual coordinate gives the
matching lower bound
$\log Z_{T+1}\ge\vartheta\sum_t\mathcal L_a(u_t)$.
Thus $T\ge1024\kappa_0^2\log(2k)$ ensures
\[
 \max_a\frac1T\sum_{t=1}^T\mathcal L_a(u_t)
 \le\frac{\log(2k)}{\vartheta T}+4\vartheta
 \le\frac1{8\kappa_0}.
\]
For $\bar u=T^{-1}\sum_tu_t$, linearity of the losses implies
\[
 \|R(B_0^{\mathsf T}\bar u-h)\|_\infty
       \le\frac{q_{\kappa,+}\rho}{8\kappa_0}.
\]
The second inequality in \eqref{eq:recovery-preconditioned-bounds}
then proves
\[
 \|\bar u\|_\infty\le\rho,\qquad
 \operatorname{opt}_{B_0}(h-B_0^{\mathsf T}\bar u)\le\rho/8\le\rho/4.
\]

Start with $h_0=g$ and $\rho_0=r_*$. At stage $j$, add the vector
$\bar u_j$ just constructed, set
$h_{j+1}=h_j-B_0^{\mathsf T}\bar u_j$, and use
$\rho_{j+1}=\rho_j/4$. Stop at
$J=\lceil\log_4(6M_0)\rceil$. The sum of stage vectors has norm at
most $\sum_{j<J}\rho_j\le4r_*/3$.
To finish exactly, choose one unit unary row at each vertex, with sign
$s_i\in\{-1,1\}$, and define the right inverse
\[
 (Ub)_{\rm unary(i)}=s_i b_i/(L_{\rm unary(i)}q_i),\qquad
 (Ub)_e=0\text{ on the other rows}.
\]
Then $B_0^{\mathsf T}Ub=b$ and
$\|Ub\|_\infty\le M_0\operatorname{opt}_{B_0}(b)$: test any feasible
$B_0^{\mathsf T}v=b$ at vertex $i$ to prove the latter inequality.
The final addition $Uh_J$ has norm at most $M_0\rho_J\le r_*/6$.
The resulting equation is exact, and the total norm is at most
$(4/3+1/6)r_*=3r_*/2$. Minimality of $J$ also gives
$\rho_J\ge r_*/(24M_0)$. The number of stages and the iteration bound
above prove \eqref{eq:recovery-static-iterations}.
\end{proof}

This construction is the elementary residual-recursion mechanism behind
Sherman's general solver, specialized to a known radius. It avoids
relative approximation to an unknown optimum. The stronger $\rho/8$
estimate leaves half the stated $\rho/4$ reduction budget for numerical
errors. The numerical section
uses absolute error budgets at the displayed stage radii, preserving a
fixed fraction of $r_*$ for snapshot perturbation and the actual final
unary residual. That argument gives a correction of norm at most $2r_*$;
it needs no lower bound on the true imbalance.

\subsection{Complementarity from three scalar identities}

There are two anchors to keep distinct. The phase fixes its budget
$K=\lambda\Delta_{\rm start}$ and its barrier references $\widehat S$.
The displacement anchor $a$ may be reset. Its current slack and mass are
\[
 S^a=c-Aa,\qquad K_a=w^{\mathsf T}S^a.
\]
In general $K_a\ne K$. Only the displacement identity uses $S^a$;
the power barrier always uses $\widehat S$.

Write
\[
 E_r=\sum_er_e\theta(r_e),\quad E_t=\sum_et_e\theta(t_e),\quad
 E_p=\sum_i\omega_i(x_i/b_0)^p,\quad
 V=\sum_i\omega_i(1+a_0/q_i),\quad Z=\sum_e\ell_ez_e.
\]
Let $\xi_x=g_x^{\mathsf T}x$ and $\xi_q=g_q^{\mathsf T}\mathbf1$.
The scalar root equation and the two Euler tests give exactly
\begin{align}
 T&=Z+\sum_e\frac{\ell_e}{L_e}
               [r_e\theta(r_e)+\eta t_e\theta(t_e)],
 \label{eq:recovery-euler-root}\\
 \ell^{\mathsf T}S^a-T-K_a+W_\sigma+\eta E_t+E_p&=\xi_x,
 \label{eq:recovery-euler-x}\\
 (1-\kappa)E_r-\kappa\eta E_t-V-E_p&=\xi_q.
 \label{eq:recovery-euler-q}
\end{align}
The first line uses the scalar root; the second tests the displacement
gradient with $x$ and uses $A(y-a)=S^a-\sigma$; the third sums the
fixed-$y$ scaling gradient. These three identities are the entire
complementarity argument.

\begin{lemma}[A usable complementarity bound]
\label{lem:recovery-complementarity}
Let $\kappa\le1/4$, $\alpha_{\rm b}\le1$, and suppose the preceding
endpoint certificate holds. Put
\[
 \beta_0=b_0q_{\kappa,-},\qquad
 \gamma=\Lambda\tau_xb_0,\qquad \zeta=\Lambda\delta_s.
\]
Assume $\zeta\le1/4$, $\gamma\le\beta_0/2$, and $\beta_0\ge16\eta$.
If the sublevel bound gives
$W_\sigma\le D$ and $\sum_e\Psi_e\le D$, and the endpoint and
anchor analyses give $V\le V_{\rm cap}$ and $E_p\le E_{\rm cap}$,
then, for
$A_*=Bm+V_{\rm cap}+E_{\rm cap}+(1+\eta)m$,
\begin{equation}\label{eq:recovery-complementarity}
 T\le4A_*+2\alpha_{\rm b}D+
                    \frac{16\eta}{b_0q_{\kappa,-}}K_a.
\end{equation}
\end{lemma}

\begin{proof}
The norm certificate and \eqref{eq:recovery-huber-elementary} give
\[
 |\xi_x|\le\gamma(E_t+m/4),\qquad
 |\xi_q|\le\zeta(E_r+m/4).
\]
Absorb the second error in \eqref{eq:recovery-euler-q}; since
$1-\kappa-\zeta\ge1/2$, the first Euler identity gives
\begin{equation}\label{eq:recovery-t-first}
 T\le Z+2(V+E_p)+2\eta E_t+\zeta m/2.
\end{equation}
Next, $E_t\le\sum_e|t_e|\le
\sum_eL_e(S^a_e+\sigma_e)/\beta_0$. In this expression the current
baseline gap cancels exactly:
\[
 \sum_eL_e(S^a_e+\sigma_e)
 =\ell^{\mathsf T}S^a+K_a+T+W_\sigma
 =2T+2K_a-\eta E_t-E_p+\xi_x.
\]
Hence
\[
 E_t\le\frac{2T+2K_a+\gamma m/4}{\beta_0-\gamma}.
\]
Substitute into \eqref{eq:recovery-t-first}. Its new coefficient of
$T$ is at most $8\eta/\beta_0\le1/2$. Absorbing it yields
\[
 T\le2Z+4(V+E_p)+\zeta m+
                    \frac{8\eta}{\beta_0}(2K_a+\gamma m/4).
\]
Finally, the scalar barrier satisfies
\[
 \ell_ez_e=B+B(\widehat S_e/z_e)^{\alpha_{\rm b}},\qquad
 \Psi_e\ge\beta_e(z_e)-(1+\eta)/2.
\]
The second inequality follows from
$a\theta(a)-\phi(a)\le1/2$ and $0<\ell_e/L_e\le1$.
Because $z_e\le\sigma_e\le S_{\max}$, the logarithmic part of
$\beta_e(z_e)$ is nonnegative. Therefore
\[
 Z\le Bm+\alpha_{\rm b}[D+(1+\eta)m/2].
\]
Substitute this bound and use $\gamma/\beta_0\le1/2$ and
$\zeta\le1/4$. All remaining terms proportional to $m$ are bounded
by $2(1+\eta)m$, giving \eqref{eq:recovery-complementarity}.
\end{proof}

The previously proved endpoint floor gives $V_{\rm cap}=O(M)$.
The anchor rule must bound $E_p=pE(x)$; a bound on $E(x)$ alone loses
the factor $p$. Keeping the barrier references fixed is equally
essential to the displayed bound for $Z$.

\begin{theorem}[Phase contraction with explicit dependencies]
\label{thm:recovery-phase-contraction}
Assume the sublevel, endpoint-floor, and anchor bounds of
Lemma~\ref{lem:recovery-complementarity}, with $0<\kappa\le1/4$ and
$0<\alpha_{\rm b}\le1$.
Suppose $D,K_a\le C_aK$, where $C_a\ge1$ is the fixed constant from
the sublevel and anchor analyses. Let $0<\varepsilon\le1/4$.
It suffices for recovery to require
\begin{equation}\label{eq:recovery-clean-tolerances}
 \tau_s\le\kappa/64,\qquad
 \tau_x\le\min\left\{\frac{\kappa}{64L_x},
                  \frac{\varepsilon q_{\kappa,-}}{8\Lambda}\right\},
 \qquad
 b_0q_{\kappa,-}\ge\frac{8(\eta+2)}\varepsilon.
\end{equation}
Then constructive recovery gives an exactly feasible nonnegative flow
whose certified gap, divided by the gap at the start of the phase, is
at most
\begin{equation}\label{eq:recovery-contraction-factor}
 \frac{4A_*}{K}
 +C_a\left(2\alpha_{\rm b}
       +\frac{16\eta}{b_0q_{\kappa,-}}+2\varepsilon\right).
\end{equation}
In particular, taking $K\ge100A_*$ and
$\alpha_{\rm b},\varepsilon\le1/(100C_a)$ makes this factor less
than $1/4$. All required inverse tolerances are polynomial in
$b_0,X,\kappa^{-1},\eta^{-1},q_{\kappa,-}^{-1}$ for constant
$\varepsilon$ and $C_a$.
\end{theorem}

\begin{proof}
The tolerances give $\delta_s\le\kappa/32$, $\zeta\le1/4$, and
$\gamma\le\varepsilon\beta_0/8\le\beta_0/32$.
They also imply $\beta_0\ge16\eta$. The correction radius in
\eqref{eq:recovery-radius} is at most $\varepsilon/4$; the known-radius construction returns a correction of norm at most
$3r_*/2<\varepsilon$. Even the numerical allowance $2r_*$ is less
than $\varepsilon$. Use that $\varepsilon$ in
\eqref{eq:recovery-flow-formula}. Combine
\eqref{eq:recovery-gap-formula} with
\eqref{eq:recovery-complementarity}, $W_\sigma\le D$, and
$D,K_a\le C_aK$, and divide by $K/\lambda=\Delta_{\rm start}$.
This gives \eqref{eq:recovery-contraction-factor}.
For the stated parameter choice, its first term is at most $4/100$.
The recovery condition implies
$16\eta/(b_0q_{\kappa,-})\le2\varepsilon$, so the remaining terms
are at most $6/100$. The sum is below $1/4$.
\end{proof}

The theorem isolates the purpose of each parameter. The phase budget
dominates the vertex and displacement energies; the small power exponent
makes the root-barrier contribution a small fraction of that budget;
the large displacement scale makes the remaining Huber contribution
small. Stationarity controls the correction and pays for its exact
construction through a separate static solve.


% Replacement dynamic-maintenance section; uses the main theorem/lemma
% environments and requires amsmath, amssymb, amsthm, array.

\section{Maintaining the two gradients by ordinary cut updates}
\label{sec:maintenance}

The optimization variables are the two additive fields
$s=\log(q/a_0)$ and $x=D_q^{-1}(y-a)$. A step adds a constant to one field on
one implicit tree cut. The implementation stores neither field as an
explicitly updated vertex array. Instead, it stores their tree
representations and a cache of complete row contributions.
The following argument applies between anchor resets, with $a,w$ and
the barrier references fixed. An anchor reset refreshes every row and
vertex; its $O(M)$ evaluations are charged separately by the reset
analysis.

There are two questions to answer. First, how much can a row's gradients
change while its cache is stale? Second, how do we detect that change
without scanning the vertices or edges of the updated cut? The same
small grounded graph answers both questions.

\subsection{The local graph and the snapshots it stores}

For a two-ended row, keep its given sign and write
\[
 A_e=a_e e_u^{\mathsf T}-b_e e_v^{\mathsf T},\qquad a_e,b_e>0.
\]
Its owner is its positive endpoint $u$. At the current state, put
\[
 C_e=Q_e e_u^{\mathsf T}-R_e e_v^{\mathsf T},\quad
 Q_e=L_ea_eq_u^{1-\kappa},\quad
 R_e=L_eb_eq_vq_u^{-\kappa},\quad
 L_e=\ell_e+w_e,
\]
and recall $r_e=C_e\mathbf1$ and $t_e=C_ex/b_0$. Define
\begin{equation}
 n_e(z)=2\left[
 R_e|z_u-z_v|+(1+|r_e|+|t_e|)|z_u|+|z_v|
 \right].
 \label{maint:localnorm}
\end{equation}
Thus one row is represented by three ordinary arcs: an edge $uv$, an
edge from $u$ to ground, and an edge from $v$ to ground. The factor two
will make the comparison with the earlier proxy exact in one direction.

For a unary row $C_e=\pm Q_e e_u^{\mathsf T}$, use
$n_e(z)=2(1+|r_e|+|t_e|)|z_u|$. This convention covers both unary signs;
there is no sign change in the physical matrix.

The vertex part of the forcing is
\[
 v_i=\omega_i\left[1+a_0/q_i+(x_i/b_0)^p\right].
\]
The \emph{true} ordinary graph norm used in this section is
\begin{equation}
 \mathcal P(z)=\sum_e n_e(z)+\sum_i v_i|z_i|.
 \label{maint:trueproxy}
\end{equation}
It has $O(m+n)$ arcs. Its frozen version
$\widehat{\mathcal P}$ has the same arc identifiers, but each row's
three capacities come from that row's last refresh, and each vertex
capacity comes from that vertex's last refresh. Different rows can
therefore have different snapshot times.

A row refresh is one synchronized operation. It queries
$s_u,s_v,x_u,x_v$ at the current state, solves that row's scalar
equation, and replaces \emph{both} coordinate-gradient contributions
and \emph{all} of that row's graph capacities. The row's motion
counters in both fields are then reset. Independent updates of the two
endpoint coefficients would not provide the error bound below.

\subsection{One local estimate controls curvature and cache error}

The following notation is confined to this proof. Suppress the row
index and set
\[
 k=q_u^\kappa\frac{\ell-w}{\ell+w}
       +\frac{\eta}{b_0}\theta(t),\qquad
 F=r\theta(r)+\eta t\theta(t).
\]
Including the row's share $-wA_ey$ of the linear objective, its two
actual coordinate gradients are the row vectors
\begin{equation}
 g_{x,e}=kC_e,\qquad
 G_{s,e}=\theta(r)C_e+kC_eD_x-\kappa F e_u^{\mathsf T}.
 \label{maint:rowgradients}
\end{equation}
In particular, $G_{s,e}$ is the derivative with $x$ fixed, not the
fixed-$y$ scaling derivative.

\begin{lemma}[Local variation of a complete row]
\label{maint:local}
Suppose the phase ranges give
\[
 |x_i|\leq X,\qquad q_i^\kappa\leq q_{\kappa,+},\qquad
 b_0\geq1,\quad 0<\kappa\leq1,\quad 0<\eta\leq1,
\]
and write
$d_0=B/(1+\alpha_{\rm b})-1-\eta>0$.
There is a universal polynomial
\[
 K_{\rm loc}=
 \operatorname{poly}(1+X+b_0+q_{\kappa,+}+d_0^{-1}+p)
\]
such that, along every smooth path in the finite-envelope domain,
\begin{align}
 |d(g_{x,e}z)|+|d(G_{s,e}z)|
   &\leq K_{\rm loc}
          [n_e(ds)+n_e(dx)]n_e(z),
             \label{maint:rowlipschitz}\\
 |dn_e(z)|
   &\leq K_{\rm loc}
          [n_e(ds)+n_e(dx)]n_e(z).
             \label{maint:normlipschitz}
\end{align}
The constants do not depend on the sizes of $Q_e,R_e,L_e$.
\end{lemma}

\begin{proof}
All derivatives below exist almost everywhere, which is sufficient for
the integrated inequalities. Write $m=n_e(ds)+n_e(dx)$.
The row split and a second, equally useful identity give
\begin{align}
 |C_ez|&\leq n_e(z),\nonumber\\
 C_eD_xz&=x_vC_ez+(C_ex-rx_v)z_u,
 \qquad
 |C_eD_xz|\leq(2X+b_0)n_e(z).
 \label{maint:multiplier}
\end{align}
We first derive the scalar transport estimate used here. Let
$\zeta(\ell)=(-\beta'_e)^{-1}(\ell)$ and write the root as
\[
 \mathcal F(\ell,r^0,t^0,\sigma)
 =\zeta(\ell)+r^0\theta((\ell+w)r^0)
       +\eta t^0\theta((\ell+w)t^0)-\sigma=0.
\]
The barrier identities $\ell\zeta\geq B$ and
$\beta''_e(\zeta)\leq(1+\alpha_{\rm b})\ell/\zeta$ give
$-\zeta'(\ell)\geq B/((1+\alpha_{\rm b})\ell^2)$. An active Huber
curvature contributes at most $1/L^2$ or $\eta/L^2$. Hence
\[
 -\mathcal F_\ell\geq d_0/\ell^2,\qquad
 |\mathcal F_{r^0}|\leq2,\qquad
 |\mathcal F_{t^0}|\leq2\eta.
\]
Implicit differentiation, using $dL=d\ell$ and $\ell\leq L$, proves
\begin{equation}
 |d\log L|\leq\frac L{d_0}
       (|d\sigma|+2|dr^0|+2\eta|dt^0|).
 \label{maint:scalartransport}
\end{equation}
Here $r^0=r/L$, $t^0=t/L$, and, with the anchor and baseline fixed,
\[
 \begin{aligned}
 Ld\sigma&=-q_u^\kappa C_e(x\odot ds+dx),\\
 Ldr^0&=C_e ds-\kappa r\,ds_u,\\
 Ldt^0&=C_e(x\odot ds+dx)/b_0-\kappa t\,ds_u.
 \end{aligned}
\]
Equations \eqref{maint:multiplier}--\eqref{maint:scalartransport}
therefore imply $|d\log L|\leq K_1m$, with $K_1$ a polynomial in the
displayed range parameters.

The remaining derivative calculation has three parts. The table states
the bounds needed in each part; $K_2,K_3$ denote universal polynomials
in the same parameters.
\[
\begin{array}{c|l}
 \text{quantity}&\text{bound}\\ \hline
 r,t&
 |dr|\leq[1+K_2|r|]m,\quad
 |dt|\leq[K_2+K_2|t|]m\\[2pt]
 R,\ 1+|r|+|t|&
 |d\log R|\leq K_2m,\quad
 |d(1+|r|+|t|)|\leq K_2(1+|r|+|t|)m\\[2pt]
 k,\ \theta(r),\ \kappa F&
 |k|\leq K_3,\quad |dk|+|d\theta(r)|\leq K_3m,\quad
 |d(\kappa F)z_u|\leq K_3m\,n_e(z).
\end{array}
\]
For the first line use
\[
 dr=C_e ds+r(d\log L-\kappa ds_u),\qquad
 dt=C_e(x\odot ds+dx)/b_0+t(d\log L-\kappa ds_u).
\]
The second line follows from the first and
$d\log R=d\log L+ds_v-\kappa ds_u$.
For the last line, $|(\ell-w)/L|\leq1$ and
$|d((\ell-w)/L)|\leq2|d\log L|$. Moreover, $d\theta(t)$ vanishes
outside $|t|\leq1$, and similarly for $d\theta(r)$. Finally,
$|d(z\theta(z))|\leq2|dz|$, and the ground coefficient in
\eqref{maint:localnorm} pays the resulting $1+|r|+|t|$ factor.

It remains to differentiate the two row vectors. Put
$\beta=d\log L+(1-\kappa)ds_u$. The exact identities are
\begin{align}
 dC_e&=\beta C_e+R(ds_u-ds_v)e_v^{\mathsf T},\nonumber\\
 d(C_eD_x)&=\beta C_eD_x+dx_uC_e
   +R[x_v(ds_u-ds_v)+(dx_u-dx_v)]e_v^{\mathsf T}.
 \label{maint:rowtransport}
\end{align}
Their last terms are weighted by $R$, the ordinary-edge capacity in
$n_e$. Thus both derivatives, tested against $z$, are bounded by a
polynomial times $m n_e(z)$. Inserting the table and
\eqref{maint:rowtransport} into \eqref{maint:rowgradients} proves
\eqref{maint:rowlipschitz}. The second line of the table proves
\eqref{maint:normlipschitz}. For unary rows the $R$ terms are absent,
so the proof is shorter.
\end{proof}

The vertex terms satisfy the same type of estimate. In actual
coordinates they are
\[
 G_{s,i}^{\rm vert}=-\omega_i(1+a_0/q_i),\qquad
 g_{x,i}^{\rm vert}=\frac{\omega_i}{b_0}(x_i/b_0)^{p-1}.
\]
Direct differentiation gives
\begin{equation}
 \begin{aligned}
 |dG_{s,i}^{\rm vert}|&\leq v_i|ds_i|,\\
 |dg_{x,i}^{\rm vert}|&\leq p v_i|dx_i|,\\
 |dv_i|&\leq p v_i(|ds_i|+|dx_i|).
 \end{aligned}
 \label{maint:vertexvariation}
\end{equation}
Here we used $|z|^{p-2},|z|^{p-1}\leq1+|z|^p$ and $b_0\geq1$.

\subsection{Global bounds, including the proxy factors}

Recall the earlier proxy $N(z)=\rho(z)+h^{\mathsf T}|z|$, where
\[
 \rho(z)=\sum_e[R_e|z_u-z_v|+|r_e||z_u|],\qquad
 h_i=v_i+\kappa\sum_{e:u_e=i}F_e.
\]
At every state, independently of stationarity,
\begin{equation}
 N(z)\leq\mathcal P(z)\leq A_N N(z),\qquad
 A_N=8+4/(\kappa\eta).
 \label{maint:proxycomparison}
\end{equation}
For the first inequality use
$\kappa r\theta(r)\leq|r|$ and
$\kappa\eta t\theta(t)\leq|t|$; the factor two in $n_e$ pays the
second copy of the row-imbalance weight. For the other inequality,
\[
 \sum_{e:u_e=i}|t_e|
 \leq\frac{h_i}{\kappa\eta}+\frac{\deg_A(i)}4.
\]
This follows from $|t|\leq t\theta(t)+1/4$. All remaining unit endpoint
weights are paid by $v_i\geq\omega_i\geq64\deg_A(i)+1$.

Summing Lemma~\ref{maint:local} and
\eqref{maint:vertexvariation} yields, after increasing $K_{\rm loc}$,
\begin{align}
 |d(g_xz)|+|d(G_sz)|
   &\leq K_{\rm loc}
      [\mathcal P(ds)+\mathcal P(dx)]\mathcal P(z),
                  \label{maint:globalvariation}\\
 |d\mathcal P(z)|
   &\leq K_{\rm loc}
      [\mathcal P(ds)+\mathcal P(dx)]\mathcal P(z).
                  \label{maint:globalnormvariation}
\end{align}
For the vertex part use $v_i\geq1$, so that
$v_i|a_i||z_i|\leq(v_i|a_i|)(v_i|z_i|)$.
In particular, along a step in either coordinate with direction $z$,
\begin{equation}
 \left|\frac{d^2}{d\nu^2}\Phi\right|
       \leq K_{\rm loc}\mathcal P(z)^2
 \label{maint:curvature}
\end{equation}
almost everywhere. This is the curvature estimate used by cut descent.
It does not assume the endpoint reverse norm comparison.

\subsection{Why frozen motion is enough}

Suppose row $e$ was refreshed at time $a$. Measure its subsequent motion
using its \emph{frozen} norm:
\[
 \mu_e^0(b)=\int_a^b[n_e^a(ds)+n_e^a(dx)].
\]
If $K_{\rm loc}\mu_e^0(b)\leq1/4$, then
\begin{equation}
 \frac12 n_e^b(z)\leq n_e^a(z)\leq2n_e^b(z),\qquad
 |(g_{x,e}^b-g_{x,e}^a)z|
 +|(G_{s,e}^b-G_{s,e}^a)z|
 \leq4K_{\rm loc}\mu_e^0(b)n_e^b(z).
 \label{maint:frozentransport}
\end{equation}
To prove this, let $\mu_e$ denote motion in the changing norm. Equation
\eqref{maint:normlipschitz} gives
$n_e^t(z)\leq e^{K_{\rm loc}\mu_e(t)}n_e^a(z)$, hence
\[
 d\mu_e\leq e^{K_{\rm loc}\mu_e}d\mu_e^0,
 \qquad
 \mu_e\leq-\frac{\log(1-K_{\rm loc}\mu_e^0)}{K_{\rm loc}}
       \leq2\mu_e^0.
\]
Integrating the gradient estimate proves
\eqref{maint:frozentransport}. The same argument applied to
\eqref{maint:vertexvariation} controls a vertex snapshot after
coordinate motion $\int(|ds_i|+|dx_i|)\leq\tau$.

Consequently, if every unrefreshed row and vertex has frozen motion at
most $\tau$, with $K_{\rm loc}\tau\leq1/4$, then
\begin{equation}
 \tfrac12\mathcal P\leq\widehat{\mathcal P}\leq2\mathcal P,
 \qquad
 |(\widehat g_x-g_x)z|+|(\widehat G_s-G_s)z|
       \leq C K_{\rm loc}\tau\mathcal P(z).
 \label{maint:cachecertificate}
\end{equation}
This statement allows different snapshot times for every row and
vertex. Each term is compared with its own current contribution before
the terms are summed.

There are two useful readings of this certificate. Relative to the
maintained graph the gradient error is at most
$2C K_{\rm loc}\tau\widehat{\mathcal P}(z)$. Relative to the earlier
proxy $N$ it is at most $C K_{\rm loc}\tau A_NN(z)$. Thus a target
error $\epsilon$ in the $N$-dual norm is obtained by choosing
\begin{equation}
 \tau\leq\min\left\{\frac1{16K_{\rm loc}},
              \frac{\epsilon}{C'K_{\rm loc}A_N}\right\}.
 \label{maint:threshold}
\end{equation}
The proxy conversion factor must be included in this choice.

\subsection{Two ordinary structures and their motion counters}

Use two independent copies of the published ordinary cut structure.
The first stores additive field $s$ and cached gradient $\widehat G_s$;
the second stores field $x$ and cached gradient $\widehat g_x$. Both
contain the same frozen graph $\widehat{\mathcal P}$, but their tree
sparsifiers and implicit cuts may differ. A step uses a cut from one
copy, updates only that copy's field, and notifies only its motion
detector. A refresh queries both copies and changes the affected graph
arcs and gradient entries in both copies. There is no transfer of an
arbitrary cut to an unrelated tree.

Both fields are pinned at ground. Equivalently, a field query returns
the stored potential at $i$ minus its stored potential at ground. If a
returned cut contains ground, use its complement and reverse the signed
step. This leaves the prescribed physical update unchanged and removes
the ordinary potential representation's irrelevant additive constant.

The ordinary interfaces are those of the tree-cut sparsifier,
directed-layer graph, and accumulated-motion detector in
van den Brand et al., arXiv:2407.10830v1, Definitions 4.2, 5.2, 5.7 and
Theorem 5.8. The dependency graph gives a superset of the edges crossing
the selected tree cut, of capacity at most $\chi$ times its true
ordinary cut capacity. The detector reports all counters reaching their
threshold, and reports only after an overestimated counter has reached
the threshold divided by its approximation factor $\Gamma$.

For each of a row's three arcs, keep one counter in each field. If an
arc has frozen capacity $c$, its counter measures $c$ times accumulated
absolute endpoint-difference motion; its threshold is $\tau/6$.
Equivalently, the significance parameter passed to the published
unweighted-motion interface is $\tau/(6c)$. All six counters share the
row's snapshot time. If any counter reports, refresh that complete row
and reset all six. A row that does not report has $\mu_e^0\leq\tau$,
as required by \eqref{maint:cachecertificate}.

For vertex $i$, its forcing arc has capacity $\widehat v_i\geq1$.
Use the significance threshold $\tau/(2p)$ in each field, measuring the
unweighted motions of $s_i$ and $x_i$. A report refreshes both vertex
gradient entries and $v_i$, at a single current state, and resets both
counters. If neither counter reports, their total true coordinate
motion is at most $\tau/p$. Thus the factor $p$ in
\eqref{maint:vertexvariation} is absorbed directly in the vertex
threshold. These are unit coordinate-motion counters; they can be
implemented on the forcing arc without adding another graph arc.

Process reports and finish their coupled refreshes before using another
cached gradient. For the line search, cap each step's normalized
length $|\nu_j|u_j$ by $\tau/(8p)$. Before a
report is processed, its old interval can then exceed the cache radius
by at most this increment. In particular, the smaller vertex radius
$\tau/p$ cannot be crossed by a step of size comparable to $\tau$.
The same frozen-transport argument, with a constant margin, remains
valid throughout the step.

\subsection{Report count and total work}

Let update $j$ have signed size $\nu_j$ and cut $S_j$, in whichever
field owns that step. Let $u_j$ be its certified upper bound on
$\widehat{\mathcal P}(\mathbf1_{S_j})$, and define
\begin{equation}
 \mathcal V=\sum_j|\nu_j|u_j.
 \label{maint:totalmotion}
\end{equation}
The total weighted counter increment at update $j$ is at most
$\chi|\nu_j|u_j$. Every row report consumes at least
$\tau/(6\Gamma)$ of this increment. Every vertex report consumes at
least $\widehat v_i\tau/(2p\Gamma)\geq\tau/(2p\Gamma)$.
Resets performed early, including the other five counters of a reported
row, only discard accumulated motion. Charging each report to the
motion it consumes therefore gives
\begin{equation}
 D_{\rm reports}=O(\Gamma\chi p\mathcal V/\tau).
 \label{maint:reportbound}
\end{equation}
This is a bound by normalized movement, not by the number of vertices
in an update. In particular, it covers edges internal to a scaling cut
because their owner or endpoint arcs cross ground.

Each report causes $O(1)$ ordinary arc changes, field queries, detector
resets, and paired vertex-gradient changes. Ground receives the negative
sum of physical gradient changes, preserving the zero-sum convention of
the ordinary interface. The two independent structures incur only a
constant replication cost. If capacities span more than a polynomial
range, use the published capacity-level decomposition; its number of
levels is polynomial in the certified logarithmic range. All these
ordinary work and approximation factors can be collected into a factor
$Q^{O(1)}\operatorname{poly}(H)$.

There is no circular maintenance assumption in the resulting bound.
For $J$ cut steps and $D_{\rm reports}$ reports, the published
output-sensitive ordinary algorithms take
\[
 (M+J+D_{\rm reports})Q^{O(1)}\operatorname{poly}(H)
\]
work, in addition to row evaluation at reports. Each row evaluation
uses its scalar root and $O(1)$ endpoint values. Equation
\eqref{maint:reportbound} bounds the reports independently by $\mathcal V$.
For several fixed-anchor intervals, add their initialization and reset
costs; this is $O(M)$ evaluations per interval.
Finally the descent estimate bounds $\mathcal V$: if a chosen cached ratio is
at least twice the error in
\eqref{maint:cachecertificate}, the true directional derivative has the
same sign and at least half its magnitude. Equations
\eqref{maint:curvature} and \eqref{maint:cachecertificate} then justify
a sufficiently small normalized step and charge its $|\nu_j|u_j$ to
potential decrease. The global phase analysis supplies the resulting
bound on $\mathcal V$.



\section{The phase algorithm and its work bound}\label{sec:descent}

\subsection{A small set of inequalities replaces the large step exponents}

Let $A_N$ be the all-state comparison factor in
\eqref{maint:proxycomparison}, and increase $K_{\rm loc}$ so it also
dominates all vertex and frozen-motion constants. Let $Q=M^{o(1)}$
dominate ordinary cut quality, detector distortion, and structural
recourse factors. A fixed power of $Q$ can be absorbed into this notation.
Choose the target endpoint norm tolerance $\epsilon_s=H^{-100}$ and set
\begin{equation}\label{eq:algorithm-scales}
 \tau=\frac{\epsilon_s}{10A_N},\qquad
 r_{\rm cache}=\frac{\tau}{100QK_{\rm loc}},\qquad
 \zeta=\frac{r_{\rm cache}}{8p}.
\end{equation}
The first converts a proxy-norm certificate into the desired $N$-norm
certificate. The second controls stale-row errors. The third prevents
one discrete step from jumping across a row or vertex cache interval.
All inverse scales are $Q^{O(1)}\poly(H)$.
Apply the maintenance estimates of Section~\ref{sec:maintenance} with
their local cache-radius parameter $\tau$ equal to $r_{\rm cache}$.
In this section $\tau$ always denotes the stopping tolerance in
\eqref{eq:algorithm-scales}.

\subsection{One cut step, with a quantitative decrease}

Extend a gradient to ground by the negative sum of its physical entries.
Its total is zero. The ordinary cut structure returns an implicit cut
$S$, its cached gradient sum $g_S\le0$, and a certified upper bound $u$
on its frozen proxy cut capacity. Its ratio approximates the largest
gradient-to-proxy ratio to within $Q$.

If $|g_S|/u>\tau/Q$, take a step of the descent sign in that field,
with dyadic size between one half and one times
\begin{equation}\label{eq:normalized-step}
 \nu=\zeta/u.
\end{equation}
Read all endpoint values needed by reports and finish all refreshes
before the next oracle use. Unreported rows have stable frozen norms;
even during the step, overshoot is at most a small fraction of the cache
radius. Thus $\mathcal P(\one_S)\le2u$ along the step.
The cached directional error is at most $\tau u/(10Q)$, and the
directional second derivative is at most $4K_{\rm loc}u^2$.
The chosen scales imply
\begin{equation}\label{eq:decrease}
 \Phi_{\rm before}-\Phi_{\rm after}
 \ge \frac{\tau}{4Q}|\nu|u
 \ge\frac{\tau\zeta}{8Q}.
\end{equation}
Indeed the linear decrease is at least $\tau|\nu|u/(2Q)$, and the
quadratic remainder is at most $2K_{\rm loc}\nu^2u^2$.
The inequality $\zeta\le\tau/(100QK_{\rm loc})$ pays that remainder.

If both structures return ratios at most $\tau/Q$, their approximation
and cache bounds give gradient dual norms at most $3\tau$ in
$\mathcal P$, and at most $3A_N\tau<\epsilon_s$ in $N$.
These are the actual coordinate-gradient certificates required in the
endpoint theorem.

The finite domain is preserved by continuation. Apply the estimates up
to the first exit from either the finite domain or the phase sublevel
and its certified range bounds. The decrease just proved prevents a
sublevel exit, and the range lemma then prevents a range exit. If a
first domain boundary contact occurred, the integrated local estimates
would keep $L_e$ and
the row metric within bounded factors of their cached values up to
that contact. At the boundary $\sigma_e=|r_e^0|+\eta|t_e^0|$, the
scalar maximizer must instead diverge. This contradiction justifies
using the local estimates throughout an admissible step.

\subsection{A reset caps displacement energy and pays for its rebuild}

Scaling steps leave $E(x)$ fixed. Along displacement steps,
\[
 |dE|\le b_0^{-1}\sum_i v_i|dx_i|\le b_0^{-1}\mathcal P(dx).
\]
Track normalized motion and scan all coordinates whenever it reaches
$M/4$ since the previous scan. Cache stability and the step cap ensure
the energy can change by at most $M$ between scans. If a scan finds
$E\ge M$, set $a\leftarrow y$, $x\leftarrow0$, and rebuild both
structures at the unchanged physical $y,q$. The barrier references,
$w$, and $\lambda$ stay fixed.

Removing the displacement Huber term from every row cannot increase
its supremum; removing $E$ decreases the potential by at least $M$.
Thus $E\le3M$ whenever an endpoint is tested and
$E_p=pE\le3pM$. Each reset costs $MQ^{O(1)}\poly(H)$ and is paid
by a decrease of $M$. A scan that does not reset is paid by the
intervening motion. Rounded resets retain at least half this decrease,
as shown in Appendix~\ref{app:numerics}.

\subsection{Parameters, ranges, and phase progress}

One admissible choice is
\begin{equation}\label{eq:parameter-choice}
 \alpha_{\rm b}=H^{-2},\quad a_0=2^{-H^8},\quad
 \kappa=H^{-9},\quad b_0=H^2,\quad K=MH^{12},\quad p=40H.
\end{equation}
The constants in $H$ can be increased once. The purposes of the choices
are easier to see than their exponents:
\begin{center}
\begin{tabular}{p{.24\linewidth}p{.66\linewidth}}
\toprule
Parameter & Required effect\\\midrule
$\alpha_{\rm b}$ & Makes $\alpha_{\rm b}D$ a small fraction of $K$.\\
$a_0,\kappa$ & Give $q_i>a_0$ at termination while keeping
$q_i^\kappa$ between constants throughout descent.\\
$p,b_0$ & Give $|x_i|=O(b_0)$ and a small weighted correction.\\
$K$ & Dominates the $O(M)$ vertex forcing and the capped $O(pM)$
displacement energy, as well as the sublevel offset.\\
\bottomrule
\end{tabular}
\end{center}
The initialization lemma with $\epsilon=H^{-4}$ and the range induction
in Appendix~\ref{sec:paired-unary-ranges} give
\[
 D=K+O(MH^8)\le2K,\quad K_a\le2K,\quad
 q_{\kappa,-}\ge1/2,\quad q_{\kappa,+}\le2,\quad |x_i|\le3b_0.
\]
The endpoint floor gives $V\le2\Omega$; the resets give
$E_p\le3pM$. The contraction factor in
\eqref{eq:recovery-contraction-factor} is then at most a constant
times $H^{-2}+H^{-11}+\varepsilon$. Take, for example,
$\varepsilon\le1/(100C_a)$ and reserve half of it for numerical
corrections. The factor is below $1/4$ after a fixed enlargement of $H$.
The recovered flow is the next baseline and the current dual vector is
the next phase anchor. Thus $O(H)$ phases reach the internal target.

\subsection{Every operation is charged}

Let $\mathcal V=\sum_j|\nu_j|u_j$ within a phase. The potential has
an initial-to-lower-bound budget $O(D)$. By \eqref{eq:decrease},
\[
 \mathcal V=O(QD/\tau),\qquad
 N_{\rm steps}=O(QD/(\tau\zeta)).
\]
The detector report estimate gives
$Q^{O(1)}p\mathcal V/r_{\rm cache}$ reports, each requiring only
$Q^{O(1)}\poly(H)$ work and constantly many graph mutations.
Scans cost $O(\mathcal V+M)$ evaluations. Resets number $O(D/M)$.
Static recovery and phase initialization have the same near-linear
bound. Since $D=M\poly(H)$ and all inverse scales in
\eqref{eq:algorithm-scales} are $Q^{O(1)}\poly(H)$, total work over
all phases is $MQ^{O(1)}\poly(H)=M^{1+o(1)}\poly(H)$.

Use deterministic ordinary cut structures. For randomized endpoint
preconditioners, choose per-phase failure at most $(M^{10}H)^{-1}$.
This is obtained by stacking $O(\log(MH))$ independent normalized
cut preconditioners: each always gives a valid congestion lower bound,
and one successful copy gives the required upper approximation.
The extra factor is $\poly(H)$, and a union bound over $O(H)$ phases
gives the stated high-probability guarantee even when $H$ is large.

\appendix


\section{Ranges, initialization, and the endpoint scaling floor}
\label{sec:paired-unary-ranges}

Let $C_0\ge\max\{1,\max_e|c_e|\}$ and
$A_{\rm row}=\max_e\|A_e\|_1$.  Take
$S_{\max}=C_0(1+A_{\rm row})$.  Recall that every coordinate has
both signed unit unary rows.  Define
\[
 q_{\max}=C_0^{1/(1-\kappa)},\qquad \Omega=\sum_i\omega_i.
\]
We use $0<\kappa\le1/2$ and $0<a_0\le1$.

\begin{lemma}[What a potential bound controls]
\label{lem:sublevel-ranges}
If $\Phi_a(y,q)\le P$, set $D=P+\Omega\log(q_{\max}/a_0)$.
Then $|y_i|<C_0$, $q_i<q_{\max}\le C_0^2$, and
\begin{equation}\label{eq:positive-budget}
 \lambda(F-d^\top y)+\sum_e\Psi_e
       +\sum_i\frac{\omega_i a_0}{q_i}+E(x)\le D.
\end{equation}
All four terms on the left are nonnegative.  Consequently,
\begin{equation}\label{eq:lower-q-range}
 q_i\ge\frac{\omega_i a_0}{D},\qquad
 |x_i|\le b_0(pD/\omega_i)^{1/p}.
\end{equation}
Let $J_D=\max\{1,\alpha_{\rm b}(D+c_\eta)/B\}$.  The row ranges are
\begin{align}
 \sigma_e-b_e&\ge\widehat S_eJ_D^{-1/\alpha_{\rm b}},
 \label{eq:finite-domain-margin}\\
 \frac B{S_{\max}}\le\ell_e
 &\le\frac{2B}{\widehat S_e}J_D^{1+1/\alpha_{\rm b}},
 \label{eq:ell-range}\\
 Z:=\sum_e\ell_ez_e&\le Bm+\alpha_{\rm b}D+\alpha_{\rm b}mc_\eta.
 \label{eq:root-moment}
\end{align}
For fixed phase data and a fixed anchor, nonempty sublevels are compact
and contained in the finite domain.
\end{lemma}

\begin{proof}
The unary slacks are $c_{i,+}-y_i$ and $c_{i,-}+y_i$.  Positivity
implies $|y_i|<C_0$.  Each exceeds $q_i^{1-\kappa}$, while their sum
is at most $2C_0$, so $q_i^{1-\kappa}<C_0$.  Every other feasible
slack is at most $|c_e|+C_0\|A_e\|_1\le S_{\max}$.
Equation~\eqref{eq:barrier-lower} therefore implies $\Psi_e\ge0$.
Primal feasibility gives
\[
 \lambda(F-d^\top y)=\sum_e w_e\sigma_e\ge0.
\]
Also $R(q)\ge\sum_i\omega_i a_0/q_i-
\Omega\log(q_{\max}/a_0)$.  These observations prove the positive
budget and the coordinate bounds.

The sandwich gives $\beta_e(\sigma_e-b_e)\le D+c_\eta$;
the root-value bound gives $\beta_e(z_e)\le D+c_\eta$.
Both arguments of $\beta_e$ are at most $S_{\max}$, so its logarithmic
term is nonnegative.  The power term yields
$\sigma_e-b_e\ge\widehat S_eJ_D^{-1/\alpha_{\rm b}}$ and
$z_e\ge\widehat S_eJ_D^{-1/\alpha_{\rm b}}$.
Substitution into \eqref{eq:barrier-inverse} proves the upper bound
on $\ell_e$; $z_e\le\sigma_e\le S_{\max}$ proves the lower bound.
Finally,
\[
 \ell_ez_e=B+B(\widehat S_e/z_e)^{\alpha_{\rm b}}
 \le B+\alpha_{\rm b}\beta_e(z_e)
 \le B+\alpha_{\rm b}\Psi_e+\alpha_{\rm b}c_\eta.
\]
Sum this inequality to obtain the bound on $Z$.
The coordinate bounds give a compact box with $q_i$ bounded away from
zero.  The row margin separates its sublevel from the finite-domain
boundary.  Continuity then makes the sublevel closed in that box.
\end{proof}

Two useful choices follow immediately.  If
$p\ge\max\{4,2\log D\}$, then $\|x\|_\infty\le e b_0$.
Certified bounds for the row scaling factors are
\[
 (a_0/D)^\kappa\le q_{u_e}^{\kappa}\le q_{\max}^{\kappa}.
\]
Compactness here proves existence of bounded sublevels, not an efficient
method for finding a stationary point.

\subsection{Starting a phase}
\label{sec:range-initialization}

\begin{lemma}[Initialization]
\label{lem:range-initialization}
Let $\widehat S_{\min}=\min_e\widehat S_e$,
$D_1=\max\{1,\max_e|A_e\one|\}$, and
$0<\epsilon\le1/[4(1+\alpha_{\rm b})]$.
If $a_0^{1-\kappa}D_1\le\epsilon\widehat S_{\min}$, then
$y=a_{\rm start}$, $q=a_0\one$, $a=a_{\rm start}$ is a finite
initial state satisfying
\[
 \Phi\le\overline P:=(1+\epsilon)K
   +B\sum_e\log\frac{S_{\max}}{\widehat S_e}
   +\frac{mB}{\alpha_{\rm b}}+4Bm\epsilon+\Omega.
\]
\end{lemma}

\begin{proof}
Here $x=t^0=0$ and
$b_e=a_0^{1-\kappa}|A_e\one|\le\epsilon\widehat S_e$.
On $[(1-\epsilon)\widehat S_e,\widehat S_e]$ direct differentiation
gives $-\beta'_e(z)\le4B/\widehat S_e$; use
$(1-\epsilon)^{-1-\alpha_{\rm b}}\le e^{1/3}<2$.
Thus the sandwich implies
\[
 \Psi_e\le\beta_e(\widehat S_e)+4B\epsilon+w_eb_e.
\]
The identity $\sum_e w_e\widehat S_e=K$ gives
$\sum_e w_eb_e\le\epsilon K$.  The objective term is $K$,
$R(a_0\one)=\Omega$, $E=0$, and
$\beta_e(\widehat S_e)=B\log(S_{\max}/\widehat S_e)+B/\alpha_{\rm b}$.
Summing proves the claim.  In particular no bound on $\max_e w_e$
was needed to initialize the potential.
\end{proof}

\subsection{Choosing the small parameter without circularity}
\label{sec:endpoint-scaling-floor}

The sublevel bound allows $a_0/q_i$ to be large during descent.  A
stationary endpoint rules this out.  The proof only needs the following
simple parameter test, rather than a separate long range bootstrap.

\begin{lemma}[Endpoint floor and a short parameter test]
\label{lem:endpoint-scaling-floor}
Use the initialization bound $\overline P$ above and put
\[
 D=D_0+\Omega\log(1/a_0),\qquad
 D_0=\overline P+\Omega\log q_{\max},\qquad
 D_{\rm col}=\max_i\sum_e|A_{ei}|.
\]
Suppose the initialization test holds and
\begin{equation}\label{eq:small-a0-test}
 \frac{D_{\rm col}(K+2B)}{\widehat S_{\min}}
       a_0^{1-\kappa}(2D)^{2+1/\alpha_{\rm b}}\le\frac14.
\end{equation}
At an endpoint with potential at most $\overline P$, assume
$g_q=C^\top\theta(r)-h$, $h_i\ge1$, and
\[
 |g_q^\top z|\le\delta_q N(z),\qquad
 N(z)\le K_N\|Cz\|_1,\qquad \delta_qK_N\le1.
\]
Then $q_i>a_0$ for every $i$, and
$\sum_i\omega_i(1+a_0/q_i)<2\Omega$.
An admissible dyadic $a_0$ exists with
\begin{equation}\label{eq:log-a0-size}
 \log(1/a_0)=O\!\left((1+\alpha_{\rm b}^{-1})
    \log\bigl((1+\alpha_{\rm b}^{-1})R_0\bigr)\right),
\end{equation}
where
\[
 R_0=\max\left\{4,D_0+\Omega+1,
  \frac{D_{\rm col}(K+2B)}{\widehat S_{\min}},
  \frac{D_1}{\epsilon\widehat S_{\min}}\right\}.
\]
\end{lemma}

\begin{proof}
All nonincreasing updates and anchor resets retain the positive budget
$D$.  The phase-start identity $\sum_e w_e\widehat S_e=K$ gives
$w_e\le K/\widehat S_{\min}$, even after a reset, since $w$ and
$\widehat S$ stay fixed.  The root ranges imply
\[
 L_{\max}\le\frac{K+2B}{\widehat S_{\min}}
                      (2D)^{1+1/\alpha_{\rm b}},
 \qquad q_j\ge a_0/D.
\]
If $q_i\le a_0$, its absolute column sum is consequently at most
\[
 \|Ce_i\|_1
 \le a_0^{1-\kappa}D^\kappa D_{\rm col}L_{\max}\le1/4.
\]
But the endpoint conditions give
$h_i\le\|Ce_i\|_1+|g_{q,i}|
\le(1+\delta_qK_N)\|Ce_i\|_1\le1/2$, contradicting $h_i\ge1$.
This proves the floor and its sum bound.  The norm comparison used
here is the endpoint certificate proved later; its proof does not use
this floor, so there is no circularity.

For an explicit choice, let $d=2+\alpha_{\rm b}^{-1}$,
$U_0=16dR_0$, and
$a_0=2^{-\lceil16d\log_2U_0\rceil}$.
Then $a_0\le U_0^{-16d}$ and
$D\le R_0(1+16d\log U_0+\log2)\le2U_0^2$.
The initialization test follows from $\kappa\le1/2$ and
$R_0\ge D_1/(\epsilon\widehat S_{\min})$.
The left side of \eqref{eq:small-a0-test} is at most
\[
 R_0U_0^{-8d}(4U_0^2)^d\le U_0^{1-5d}<1/4.
\]
This proves existence and \eqref{eq:log-a0-size}.
\end{proof}

In particular, a bound on $\log(1/\widehat S_{\min})$ yields a bound
on the working ranges.  No bound on input denominator lengths is
implicit in this statement.  Propagating the reference-slack bound
between phases remains a separate induction.

\subsection{Checking the displayed parameter choice across phases}
\label{app:parameter-check}

Here we verify \eqref{eq:parameter-choice}, rather than treating its
range requirements as extra hypotheses. The initial padded data have
logarithmic magnitude $O(H)$ and reference slacks at least one.
The initial primal--dual gap is at most a polynomial in the original
data magnitudes and $M$. Set the internal accuracy in advance and let
$N_{\max}$ be the number of factor-four contractions needed to reach it.
Then $N_{\max}=O(H)$ is known before the phases run.

For the first $j\le N_{\max}$ phases maintain inductively
\begin{equation}\label{eq:reference-induction}
 \log(1/\widehat S_{\min})\le CjH^3
\end{equation}
for a universal constant $C$. This holds initially. In the current
phase, $\log S_{\max}=O(H)$, $\Omega=O(M)$ and
$\log q_{\max}=O(H)$. Use initialization tolerance $\epsilon=H^{-4}$.
Its potential bound gives
\[
 \overline P=K+O(MH^8),\qquad
 D=\overline P+\Omega\log(q_{\max}/a_0)
      =K+O(MH^8)\le2K.
\]
The term $\epsilon K=MH^8$ is the largest new initialization offset;
the negative logarithmic regularizer has the same order. The barrier
logarithms cost only $O(MH^4)$, and its power offset costs $O(MH^2)$.
All constants are absorbed by a fixed enlargement of $H$.

Since $\log D=O(H)$, the root bounds imply
\[
 \log L_{\max}=O(H^4),\qquad
 \log(1/L_{\min})=O(H),\qquad
 \log(1/(\sigma_e-b_e))
 \le\log(1/\widehat S_e)+O(H^3).
\]
Rounded anchor changes can consume at most a factor two of this slack
margin. The next phase's reference slack therefore obeys
\eqref{eq:reference-induction}, after increasing $C$. Keeping references
fixed during within-phase resets is what permits this induction.

Initialization and the endpoint floor test also follow quantitatively.
The logarithm of their left sides contains the negative term
$-(1-\kappa)H^8\log2$. All positive terms total at most $O(H^4)$:
they come from $\widehat S_{\min}^{-1}$, $D_{\rm col}$, $K$, and
$(2D)^{2+1/\alpha_{\rm b}}$, whose logarithm is $O(H^3)$.
Thus both inequalities hold with large fixed slack. The endpoint norm
certificate used in the floor test has no dependence on this floor,
so the order of proof is not circular.

Throughout descent, the coarse lower bound is sufficient:
\[
 q_i\ge a_0/D,\qquad
 q_i^\kappa\ge\exp[-\kappa(H^8\log2+\log D)]\ge1/2,
 \qquad q_i^\kappa\le q_{\max}^\kappa\le2.
\]
Also $p=40H$ gives $(pD)^{1/p}\le3$, so $|x_i|\le3b_0$.
Only at a stationary endpoint do we use the stronger $q_i>a_0$.
It gives $V\le2\Omega=O(M)$, whereas the energy cap gives
$E_p\le3pM=O(MH)$. Substituting in the contraction theorem proves
the required factor-four progress for the current phase. This completes
the induction for all $N_{\max}$ phases and simultaneously justifies
the phase count. In particular no bound involving $1/(1-\gamma_e)$
has entered the argument.



\section{From original capacities to the grounded program}\label{app:reduction}

For an original arc $e=(u,v)$, introduce a new coordinate $j_e$ and two
nonnegative variables $f_e,z_e$, with rows
\[
 A_{f_e}=-e_u^{\mathsf T}+\gamma_e e_{j_e}^{\mathsf T},\qquad
 A_{z_e}=-e_v^{\mathsf T}+e_{j_e}^{\mathsf T}.
\]
Give them costs $c_e$ and zero. The new demands are
\[
 d'_{j_e}=\gamma_eu_e,\qquad
 d'_v=d_v-\sum_{e\text{ entering }v}\gamma_eu_e.
\]
The equation at $j_e$ reads $\gamma_ef_e+z_e=\gamma_eu_e$.
It is exactly the capacity condition $f_e\le u_e$, with
$z_e=\gamma_e(u_e-f_e)$. At an old vertex the new row sum is
$-\sum_{\rm out}f_e-\sum_{\rm in}z_e$; substituting for $z$ gives
the original net-inflow equation. Thus this capacity gadget is exact
and keeps at most two nonzeros of opposite signs in each row.

Now add both signed unit unary variables at every coordinate, at common
positive cost $U$. They ensure an easy primal start, a bounded dual
region, and an exact residual repair. For target original conservation
error $\epsilon_{\rm out}$, one sufficient choice is
\[
 U=4W^2+
       \frac{8(1+W)^2(2m_{\rm original}W^2+1)}{\epsilon_{\rm out}}.
\]
At old vertices use dual anchor zero and at gadget coordinates use
$-2W^2$. Every physical-row slack is at least one; the unary costs
also leave strict slack. A feasible padded primal baseline uses
$f_e=0,z_e=\gamma_eu_e$ and signed unary flow for the remaining original
demands. All initial magnitudes and inverse positive slacks have
logarithm $O(H)$ when $\log(1/\epsilon_{\rm out})=O(H)$.

The artificial rows generally do change exact feasibility after they
are removed. The next estimate is the precise reason they are allowed.

\begin{lemma}[Removing padding and clipping]\label{lem:padding}
Suppose the original instance is feasible. Let a feasible padded solution
have artificial mass $a$ and objective at most $\OPT+\epsilon$.
If $U>W^2$, clipping its original coordinates to $[0,u]$ gives
\[
 c^{\mathsf T}\bar f\le\OPT+\epsilon,\qquad
 \norm{B^{\mathsf T}\bar f-d}_\infty
 \le [2+W(1+W)]\frac{2m_{\rm original}W^2+\epsilon}{U-W^2}.
\]
\end{lemma}
\begin{proof}
Write the signed unary contribution as $r$, with $\norm r_1\le a$.
The gadget equation gives $f_e\le u_e+W|r_{j_e}|$, hence
$\norm{f-\bar f}_1\le Wa$. Substitution into the old equations gives
imbalance $-r_v-\sum_{e\text{ entering }v}r_{j_e}$. Its $\ell_1$ norm
is at most $2a$. Clipping changes that norm by at most
$(1+W)\norm{f-\bar f}_1\le W(1+W)a$ and changes objective by at most
$W\norm{f-\bar f}_1\le W^2a$.

Since $f_e\le u_e+W|r_{j_e}|$ and $|c_e|\le W$,
$c^{\mathsf T}f\ge-m_{\rm original}W^2-W^2a$.
An original feasible optimum has absolute objective at most
$m_{\rm original}W^2$ and is a feasible padded solution with no
artificial mass. Consequently
$(U-W^2)a\le2m_{\rm original}W^2+\epsilon$.
Finally,
$c^{\mathsf T}\bar f\le c^{\mathsf T}f+W^2a
\le\OPT+\epsilon-(U-W^2)a$, proving both bounds.
\end{proof}

For maximum delivery, delete the source and sink columns from the
original balance equations, set the remaining demands to zero, and use
cost $-B_{et}$ on original arc $e$. Missing endpoints can make a row
unary; its actual coefficient and sign are retained. The same gadget
and padding argument then give an approximate maximum-delivery flow.
Use the stricter conservation accuracy specified in the next appendix.
The corresponding $U$ still has logarithm $O(H)$.



\section{Exact internal conservation for lossy maximum delivery}
\label{app:lossy-repair}

In this appendix $m,n$ refer to the original retained graph. Here $0<\gamma_e\le1$. The output must have exact balance at all
vertices other than source $s$ and sink $t$. Approximate padded
feasibility alone does not imply this guarantee.

\subsection{Why a normalized acyclic optimum exists}

Start with an optimal exactly feasible original flow. If its positive
support has a directed cycle, subtract the same amount
$\varepsilon=\min_{e\text{ on cycle}}f_e$ from every cycle edge.
At an internal cycle vertex, net inflow increases by
$(1-\gamma_{\rm incoming})\varepsilon\ge0$.
Temporarily allow this nonnegative excess. If the cycle contains $t$,
its net inflow also increases; otherwise its objective is unchanged.
Each cancellation removes an edge, so after finitely many cancellations
the support is a DAG with nonnegative internal excess and no smaller
sink value.

Such a DAG flow is an existence-only sum of generalized paths. A path
starts at $s$ or $t$ and ends at $s,t$, or an internal vertex with excess.
This decomposition follows by tracing positive outgoing flow backward:
an internal vertex cannot be a net source, and acyclicity ends the trace.
Flow along a path is attenuated by the successive gains. Keep only
paths from $s$ to $t$. A discarded path starting at $t$ contributed
negatively to sink net inflow; all other discarded paths contributed
zero. The kept flow has exact internal balance, no incoming source
arcs, no outgoing sink arcs, and no smaller delivery. It is therefore
also optimal. Its exact DAG support decomposes into at most $m$ paths,
because maximal subtraction of a positive $s$--$t$ path kills an edge.
No path decomposition is executed by the algorithm.

We may therefore delete incoming arcs at $s$, outgoing arcs at $t$,
and vertices unable to reach $t$. If no arc remains, return zero. Otherwise put $F_{\max}=m u_{\max}$.
If $F_{\max}\le\delta$, the zero flow is sufficient.

\subsection{Pruning with a certified source tree}

Let $a_*=\delta/(8F_{\max})$. We need source-tree paths of noticeable
gain, without distinguishing nearly equal rational path gains.
Use multiplicative Dijkstra with a monotone downward rounding map
\[
 F_e(g)=\operatorname{round}_{\downarrow}(\gamma_e g),
 \qquad (1-\xi)\gamma_e g\le F_e(g)\le\gamma_e g,
 \quad \xi=1/(10n),
\]
above the truncation threshold $a_*/4$; smaller labels are set to zero.
The map is monotone and never exceeds $g$, so the usual label-setting
proof applies. Parent pointers always point to previously settled
vertices and form simple paths. The nonzero working labels have
logarithmic range $O(H)$ and require $O(\log n)$ relative mantissa bits.

Retain vertices whose settled label is at least $a_*$. Their true
parent-path gain is at least that label. If a vertex's true maximum
source gain is at least $2a_*$, take a simple maximizing path.
Every prefix has at least its final gain, because all gains are at most
one. Rounding along at most $n-1$ edges loses a factor no smaller than
$(1-\xi)^{n-1}>1/2$, so that path never truncates and its endpoint
receives label at least $a_*$. Thus a deleted vertex has true maximum
source gain below $2a_*$.

\emph{Recompute reverse reachability from $t$ in the retained graph.}
Delete vertices that no longer reach $t$. This second deletion is
necessary: a retained vertex's former route to $t$ might have passed
through a low-gain vertex. It removes no surviving $s$--$t$ path.
It also preserves source-tree prefixes: a retained prefix vertex
reaches its retained descendant, and then $t$.
If no source--sink path remains, return zero; the same pruning bound
below certifies its error.

In the existence-only optimal path decomposition, a path meeting a
low-gain vertex delivers at most $2a_*u_{\max}$. Charging each of at
most $m$ paths once loses at most $2a_*F_{\max}=\delta/4$.
The final retained graph now has both a source tree of gain at least
$a_*$ and a directed tree toward $t$.

\subsection{The two sweeps and their capacity cost}

On that graph obtain a clipped approximate flow with objective error
at most $\delta/4$ and absolute internal balance errors at most
\begin{equation}\label{eq:repair-accuracy}
 \epsilon\le c\frac{\delta a_*u_{\min}}{nF_{\max}}
\end{equation}
for a sufficiently small universal $c$. Its required logarithmic
accuracy is $O(H)$, since $u_{\min}\ge1/W$.

First process the source tree from leaves to source. At an internal
vertex with deficit, add just enough flow on its incoming tree edge
to remove it, and charge the withdrawal to the parent. Stop at $s$.
Each initial deficit contributes at most its size divided by the
gain of a retained source path to any upstream edge. Therefore no
edge receives more than $n\epsilon/a_*$. A request can consume an
existing positive excess but cannot create one. At the end every
internal balance is nonnegative and the total remaining internal
excess is at most $n\epsilon$.

Second process a tree directed toward $t$ from leaves to root. Send
each internal vertex's excess on its outgoing tree edge. The received
amount at its parent is no larger. Hence no edge receives more than
$n\epsilon$, and every internal balance becomes exactly zero.
Neither sweep uses negative flow or decreases sink inflow.

The combined additions violate each capacity by at most a factor
$1+\theta$, where
$\theta\le2n\epsilon/(a_*u_{\min})$.
Divide all flows by $1+\theta$. Exact balance is homogeneous, so it
is preserved. The lost delivery is at most
$\theta F_{\max}\le\delta/4$ by \eqref{eq:repair-accuracy}.
Indeed the unscaled repaired sink inflow is at most
$(1+\theta)F_{\max}$, by its relaxed capacity bounds.
Together with pruning and optimization, the total loss is below
$\delta$. The sweeps cost $O(M)$ original-coefficient operations;
the rounded path computation costs $O(M\log M)$ operations.



\section{Numerical realization with an absolute correction budget}
\label{app:numerics}

The precision argument has two parts: show that every scale used as
a denominator has a bounded logarithmic range, and use absolute
error budgets for all cancellation quantities. It must not assume
that the true imbalance, a gradient entry, or $Q_e-R_e$ is separated
from zero.

\subsection{What is stored and what is recomputed}

Store all additive cut increments for $s$ and $x$ as exact dyadics.
Their tree representations therefore have no floating-point trajectory
error. The centered field starts at exact zero. At a row refresh,
query both fields, evaluate $q_i=a_0\exp(s_i)$
and $y_i=a_i+q_ix_i$ at the queried endpoints, solve the scalar row
equation, and then round the complete row contribution. The two
coordinate gradients and every capacity of that row use the same
snapshot. Approximate evaluations affect certified choices and cache
errors; they do not replace the stored additive fields by a perturbed
uncontrolled trajectory.

The range induction gives a universal $C$ for which positive
denominator scales, scalar-root conditioning scales, and finite-domain
margins lie in $[e^{-H^C},e^{H^C}]$. Initial data, fields, row weights,
matrix entries and gradient magnitudes are bounded above by $e^{H^C}$.
Only the first class has a lower bound. In particular $w_e$ may vanish,
and $\ell_e-w_e$, residuals, and derivative coordinates may be
arbitrarily small.

There are at most $e^{H^C}$ operations after enlarging $C$ to dominate
the proved work count and all polynomial factors. Choose common
absolute evaluation error $\epsilon_{\rm num}=2^{-H^{C'}}$, with a
fixed $C'>C+10$. This makes the sum of evaluation, dyadic-step,
and cache-rounding errors smaller than a tenth of the reserved descent,
cache, domain, and phase-gap margins. Exponentials and powers use range
reduction and series with certified remainders. Scalar roots use
logarithmic bracketing and bisection; the explicit lower bound on the
root derivative certifies the final interval. These computations use
$\poly(H)$ bits and operations.

Every ordinary graph ground coefficient is stored with its positive
floor, such as $1+|r|+|t|$ or $h_i+\sum_{e:u_e=i}|r_e|$.
An arbitrarily small $|r_e|$ is merged into that ground coefficient,
rather than treated as a separate tiny edge. Binary coefficients have
the lower bound supplied by $L_{\min},q_{\min}$ and the original
nonzero row coefficients. Thus capacity-level decomposition has
$\poly(H)$ levels. A common positive scaling of graph weights is
handled without changing normalized steps: multiply capacities and
gradients by $D_*$ and store the field divided by $D_*$.

\subsection{Why rounded static recovery retains its radius}

Write $T_B=B_0^{\mathsf T}$ for the ideal endpoint correction operator
and $U$ for its explicit unit-unary right inverse. Choose a known
absolute radius $r_*>0$ from \eqref{eq:recovery-radius}, with a constant
fraction of the permitted recovery radius still unused. This bound is
positive even when the actual correction optimum is zero.

Form a consistent recovery snapshot as follows. Choose positive rational
$\widetilde q,\widetilde\ell$ approximating the ideal values, keep
$w=\lambda f^0$ and $\lambda$ as exact phase expressions, and set
\[
 \widetilde L=\widetilde\ell+w,\qquad
 \widetilde B_0=\diag(\widetilde L)A\diag(\widetilde q),\qquad
 \widetilde g=\diag(\widetilde q)A^{\mathsf T}(\widetilde\ell-w),
 \qquad \widetilde T_B=\widetilde B_0^{\mathsf T}.
\]
The matrix $A$ retains its original coefficients. In particular $w$ is
not independently rounded, which would destroy $A^{\mathsf T}w=\lambda d$.
All snapshot entries are rational original-coefficient expressions;
no exact exponential or exact scalar root is needed in a final residual.
Choose $\widetilde q,\widetilde\ell$ so the resulting positive unary
denominators differ from their ideal values by relative error below
$1/4$. Define $\widetilde U$ using those exact snapshot denominators;
then this right inverse exists and has
norm at most twice the ideal elementary bound. For any ideal witness
$T_Bv_*=g$, $\|v_*\|_\infty\le r_*$, the explicit perturbed witness is
\begin{equation}\label{eq:perturbation-witness}
 \widetilde v_*=v_*+\widetilde U
     [\widetilde g-g-(\widetilde T_B-T_B)v_*].
\end{equation}
It satisfies $\widetilde T_B\widetilde v_*=\widetilde g$ exactly.
If the demand error is at most $a$ in $\ell_\infty$ and every matrix
entry error is at most $a$, then
\[
 \|\widetilde v_*-v_*\|_\infty
 \le\|\widetilde U\|_{\infty\to\infty}\,a(1+mr_*).
\]
The right-inverse norm is at most $e^{H^C}$. Our precision makes this
less than $r_*/100$. The snapshot graph norm and its condition
inequalities also change by only a reserved constant factor: positive
graph coefficients are rounded relatively. More explicitly, if
$\mu=\|\widetilde U\|\,\|\widetilde T_B-T_B\|\le1/100$, then the same
witness with no demand perturbation gives
$\operatorname{opt}_{\widetilde T_B}(b)
\le(1+\mu)\operatorname{opt}_{T_B}(b)$ for every $b$.
Interchanging the two operators gives the reverse inequality, using
the ideal right inverse. Thus their minimum-congestion norms are
uniformly equivalent, and the preconditioner condition changes by only
a constant factor. Use the snapshot starting radius
$r_0=(1+1/50)r_*$, rounded upward by less than $r_*/1000$ to a rational
value if necessary. The perturbed witness fits within this radius.
The exact solver's norm bound is at most $3r_0/2<1.54r_*$.
Reserve an additional $r_*/50$ for returning to the original snapshot
coefficients and for arithmetic errors. The resulting bound is below
$2r_*$. Choosing $r_*\le\varepsilon/10$ leaves ample space in the
nonnegative recovery formula.

This is an additive perturbation argument. It does not assert that
fixed absolute precision preserves a relative factor times an
arbitrarily tiny or zero correction optimum.

\subsection{The multiplicative-weights computation needs no tiny signs}

The solver uses stage radii $\rho_j=r_0/4^j$ and stops after
$O(\log M_0)$ stages. Thus every normalization radius is at least
$r_*/(24M_0)\ge e^{-H^C}$, independent of the size of the actual
imbalance. Each stage has exact residual-reduction slack, from
$\rho_j/8$ to the required $\rho_j/4$.
Use the slack for absolute errors in matrix products, weights and
support decisions. If $B_0R^{\mathsf T}(p^+-p^-)$ is approximated
coordinatewise to error $a$, any resulting cube support choice is
suboptimal by at most $2m\rho_j a$ before normalization, or
$2ma/q_{\kappa,+}$ in normalized loss. Near-zero entries need no
sign resolution.

For $T_j$ iterations, require the support-product coefficient accuracy
and the normalized loss/weight accuracy, respectively, to satisfy
\[
 a\le\frac{q_{\kappa,+}}{10^4mT_j\kappa_0},\qquad
 \epsilon_{\rm loss}\le\frac1{10^4T_j\kappa_0}.
\]
The factor $m$ explicitly pays for summing coefficient errors over the
cube coordinates. Cumulative support and loss errors then consume less
than half the stage slack.
Compute exponential weights by subtracting the largest log weight and
truncating tails with certified total mass below this accuracy.
Only $O(\log(kT_j\kappa_0))=\poly(H)$ logarithmic range is retained;
there is no requirement to exponentiate a large negative cumulative
logit to relative precision. All stage scales and inverse accuracies
have logarithm $\poly(H)$.

Finish with the unit-unary correction of the \emph{actual} snapshot residual
$\widetilde g-\widetilde B_0^{\mathsf T}v$, evaluated using the original coefficients and
the returned vector. Its norm is bounded by the coarse-solver estimate
plus the additive perturbation bound above. This corrects the actual
snapshot, rather than an earlier unrounded vector.
Apply \eqref{eq:recovery-flow-formula} with $\widetilde\ell,w,\lambda$
and this exactly corrected vector. The consistent snapshot definitions
give $A^{\mathsf T}(\widetilde\ell-\widetilde L\odot v)=\lambda d$
identically, so the recovered flow conserves the original $A$ exactly.
The complementarity numerator changes by at most
$\sum_e\sigma_e|\widetilde\ell_e-\ell_e|
\le e^{H^C}\epsilon_{\rm num}$; this is inside the reserved phase-gap
margin. Thus exact snapshot feasibility and the certified objective
bound hold for the same returned flow.

\subsection{Anchors, primal baselines, and final exact expressions}

At a reset, replace $a=y$ by a rational approximation within
$\epsilon_{\rm num}$ and set $x=0$. This moves the physical dual
vector by that amount; it is not an exact $y$-preserving operation.
The bounded derivatives and domain margin make the resulting potential
increase negligible. The reset still decreases potential by at least
$M/2$. Barrier references and $w$ remain fixed. The same construction
starts the next phase with a rational strictly feasible anchor.

Compute $w$ and the baseline objective once per phase. Test a certified
interval for the phase gap before dividing by it. An interval overlapping
the internal stopping range is handled using its certified upper end,
within the reserved output margin.

Round a recovered nonnegative primal flow down to dyadics, with error
at most $\epsilon_{\rm num}$ per entry. Repair its actual
original-coefficient balance residual with signed unit unary variables.
For a residual $r_i$ and approximation $\widehat r_i$ with error at most
$\upsilon$, use
\[
 k_i^+=\max\{0,-\widehat r_i\}+\upsilon,\qquad
 k_i^-=k_i^++r_i.
\]
Both entries are nonnegative and $k_i^+-k_i^-=-r_i$ exactly.
The added mass is $O(MW\epsilon_{\rm num}+n\upsilon)$, with negligible
objective cost. This keeps each next baseline in a fresh $O(M)$-size
shared expression representation. Only original-coefficient arithmetic
is used to represent the exact residual; its potentially tiny sign is
never decided by a floating-point comparison.

After removing artificial variables, minimum-cost output has the
approximate balance guarantee of Lemma~\ref{lem:padding}. For lossy
maximum delivery, the final sweeps use original gains and preserve
balance identically. Fully expanding those expressions can depend on
input denominator lengths, which are outside the arithmetic bound.


\begingroup\small
\begin{thebibliography}{9}\setlength{\itemsep}{3pt}
\bibitem{cuts} J. van den Brand, L. Chen, R. Kyng, Y. P. Liu,
S. Meierhans, M. Probst Gutenberg, S. Sachdeva.
\emph{Almost-Linear Time Algorithms for Decremental Graphs:
Min-Cost Flow and More via Duality}. FOCS 2024.
\url{https://arxiv.org/abs/2407.10830}.
\bibitem{sherman} J. Sherman.
\emph{Generalized Preconditioning and Network Flow Problems}.
\url{https://arxiv.org/abs/1606.07425}, version 2.
\bibitem{generalized} S. Jiang, M. Kapralov, L. Li, A. Sidford.
\emph{Generalized Flow in Nearly-linear Time on Moderately Dense Graphs}.
FOCS 2025. \url{https://arxiv.org/abs/2510.17740}.
\end{thebibliography}
\endgroup
\end{document}
